\documentclass[12pt]{article}
\title{Digest of The New York Times}
\author{The New York Times}

% % \definecolor{dark-green}{rgb}{0,0.25,0}
% \usepackage[utf8]{inputenc}

\usepackage{fontspec,xunicode,xltxtra}
\setmainfont[Ligatures=Required,Mapping=tex-text]{Georgia}
% \setmainfont[Ligatures=Required,Mapping=tex-text]{Palatino Linotype}
\setsansfont[Ligatures=Required,Mapping=tex-text]{Myriad Pro}
\setmonofont[Ligatures=Required,Mapping=tex-text]{Courier New}
\usepackage{xeCJK}
\setCJKmainfont{宋体}
\setCJKmonofont{黑体}

% \usepackage{indentfirst}


\usepackage[dvips,dvipsnames,svgnames]{xcolor}
\definecolor{light-gray}{gray}{0.95}

%生成PDF的链接
\usepackage[dvipdfmx]{hyperref}
\hypersetup{
    bookmarks=true,         % show bookmarks bar?
    bookmarksopen=true,
    pdfpagemode=UseNone,    % options: UseNode, UseThumbs, UseOutlines, FullScreen
    pdfstartview=FitH,
    pdfborder=1,
    pdfhighlight=/P,
    pdfauthor={wuxch},
    unicode=true,           % non-Latin characters in Acrobat’s bookmarks
    colorlinks,             % false: boxed links; true: colored links
    linkcolor=blue,         % color of internal links
    citecolor=blue,        % color of links to bibliography
    filecolor=magenta,      % color of file links
    urlcolor=cyan           % color of external links
}
% 使用目录
\usepackage{makeidx}
% 使用图形
\usepackage{graphicx}


\usepackage{titletoc}
\dottedcontents{section}[1.5em]{}{2.3em}{1pc}

% % 页眉页脚的设定
% \usepackage{fancyhdr}
% \pagestyle{fancy}
% \lhead{\includegraphics[height=0.2in]{zte_logo}}

% \usepackage{footmisc}
\usepackage{pfnote}
\usepackage{fnpos}

% % 嵌入的代码显示
% \usepackage{listings}
% \lstset{language=C++, breaklines, extendedchars=false}
% \lstset{keywordstyle=\color{blue},
%         commentstyle=\color{SeaGreen},
%         stringstyle=\ttfamily,
%         showstringspaces=false,
%         backgroundcolor=\color{light-gray}}
% \lstset{basicstyle=\small}

\usepackage{booktabs}

\usepackage{tipa}
\usepackage{phonetic}
\usepackage[gen]{eurosym}

% \usepackage{CJK}
\usepackage{makeidx}
% \usepackage{glossary}

\usepackage{fancyhdr}

\usepackage{lastpage}
\cfoot{\textsf{\thepage\ of \pageref*{LastPage}}}

\pagestyle{fancy}
\fancyhead{}
\renewcommand{\headrulewidth}{0pt}
\pagenumbering{arabic}

\usepackage{graphicx}
\usepackage{picinpar}

\setlength{\parindent}{0in}
\setlength{\parskip}{2.5ex plus4mm minus3mm}
% \setlength{\parskip}{.3\baselineskip}

\setcounter{section}{0}

% bibtex
\usepackage[super]{natbib}
% \usepackage[globalcitecopy,labelstoglobalaux]{bibunits}
% \makeatletter
% \renewcommand{\@bibunitname}{\jobname.\the\@bibunitauxcnt}
% \makeatother

% \renewcommand{\refname}{缩略语}
% \renewcommand{\refname}{}

\makeatletter                   % Reference list option change
\renewcommand\@biblabel[1]{#1:  } % from [1] to 1: 
\makeatother                    %

\newcommand{\citenumfont}[1]{\textsf{#1}}

% \usepackage{tipa}
% \usepackage{phonetic}

\newcommand{\phonetic}[1]{{\fontspec[Colour=CC00CC]{Microsoft Sans Serif} #1}}


%%% Local Variables:
%%% TeX-master: "medical_trivia"
%%% End:

\usepackage{config}

\makeindex
\begin{document}
\date{}
\thispagestyle{empty}

\begin{figure}
\includegraphics*[width=0.7\textwidth]{The_New_York_Times_logo.png}
\vspace{-15ex}
\end{figure}
% \renewcommand\contentsname{\textsf{Digest of The New York Times}}
\renewcommand\contentsname{}
{\footnotesize\textsf{\tableofcontents}}

\clearpage
\setcounter{page}{1}

\section{Web Is Bringing Skaters' Lives Uneasily Close to Fans' Obsessions}

\lettrine{W}{hen}\mycalendar{Jan.'10}{25} Brian Orser received an e-mail message last fall from an
admirer of his star figure skater, Kim Yu-na of South Korea, he wrote it off as another piece of fan
mail.

\myphoto{articleInline.jpg}{Social media and the Internet have made it harder for top skaters like
  South Korea's Kim Yu-na to keep obsessive fans at arms' length.}

The fan, a Canadian teenager, gave his name and attached his photograph, asking Orser for his help
in setting up a rendezvous with Kim, the gold-medal favorite at the Vancouver Olympics. Orser, Kim's
coach, wrote back and said thanks, but no.

``I quite regularly get e-mails from people sending Yu-na their support or best wishes,'' said
Orser, a two-time Olympic medalist for Canada. ``And this one seemed fine, until a few more e-mails
came from the same e-mail address, but now the guy was using a different name each time. He was
professing his undying love for her and that he would do anything to meet her.''

In figure skating, those obsessions are nothing new. To see some of the skaters who soon will
command the spotlight in Vancouver, fans this week have traveled to the United States Figure Skating
Championships here from as far away as Russia and Japan.

They are drawn to the elements of skating that often make the sport a television hit every four
years: the costumes, the music and the graceful beauty and athleticism of young men and women.

Those elements, however, seemingly conspire to create fans who take their interest to the verge of
obsession -- and beyond.

Overzealous fans may be attracted more to figure skaters than to other athletes because of the
skaters' fancy and often skimpy attire in competition, security experts say. A person may lock on to
one athlete because the skater is good looking or because the person believes he or she has a deep
connection with the skater. Those obsessions at times may be fueled by mental illness.

Some top skaters have had firsthand experiences with fans who crossed the line.

``Some of them think they are your best friend,'' said Sasha Cohen, the 2006 Olympic silver
medalist, adding that the Internet had provided fans with more information than ever, including
simple ways to track down their favorite athlete. ``They want to fly you to their school, so you
could be their show-and-tell. One grandfatherly-like type character wrote letters, then actually
flew across the country to take me to Disneyland. He showed up at my rink.''

The letters sent by fans to skaters are mostly ones of encouragement, skaters say, but there are
also some that are questionable. The ones that make athletes nervous or feel threatened are
forwarded to Larry Buendorf, the chief security officer for the United States Olympic Committee.

Buendorf, a former Secret Service agent, said he had several open files on fans who have unusual
interest in current skaters.

``Figure skating seems to attract more of these kinds of letters than other sports, but I've seen it
all,'' said Buendorf, who will be at the Olympic skating competition in Vancouver to ensure the
safety of American athletes. ``For the most part, there's no problem with having fans write letters,
saying, I want to marry you, or I think you're wonderful, but that can easily escalate to another
level. If the person starts trying to meet the athlete or shows up and takes that athlete out of
focus, that's when we would step in.''

Buendorf said he could count on one hand the number of serious incidents during his 17 years with
the U.S.O.C. One of those was in 1996, when he was asked to help protect Michelle Kwan, a five-time
world champion, from a fan who had written death threats after she had not responded to his previous
letters. That fan was taken into custody in Paris when Kwan was competing in the Grand Prix final
there.

Since then, the relationship between fan and athlete has changed. The emergence of social media and
other forms of electronic communication, like e-mail and texting, is bringing fans closer than ever
to the athletes.

``The controls are not in place anymore,'' Buendorf said. ``The fan mail is going directly to the
athlete.''

Buendorf says he encourages athletes to use the Internet with caution and post only general
information like, ``I did a great job,'' or ``What a great event.'' If they do receive a disturbing
message, he tells them not to delete it, so authorities can track its origin.

The old-fashioned way to communicate with an athlete was to send a letter to their rink, coach, fan
club or national federation. Most times, the skater never saw the ones that might have raised red
flags because those letters were screened by parents, agents or coaches. U.S.~Figure Skating
forwards fan mail to its athletes with a warning.

``We say, 'This hasn't been opened, so we encourage you to exercise caution in opening it or in
dealing with people you don't know,' '' said Scottie Bibb, a spokeswoman for the organization. ``The
letter we send is pretty thorough, especially because a lot of our skaters are underage and may not
be as cautious as they should be.''

Jody Flatt, the mother of the Olympic hopeful Rachael Flatt, says she always tries to intercept her
daughter's fan mail. She has advised Rachael to watch what she posts on her Twitter page.

``When it comes to her daily plans, like traveling or going places, we always post them after the
fact, not in advance,'' Jody Flatt said of Rachael, 17. ``We don't want people showing up where she
is, out of the blue.''

U.S.~Figure Skating encourages the same policy. Bibb said the organization also warned its skaters
to be careful whom they accepted as friends on Facebook.

``If you friend everyone on the entire planet, then say, `Hey, I'm headed to the movies at this time
and this place,' it's almost like saying, 'Come and get me,' '' Bibb said. ``They have to remember
that they are like rock stars and people want to be close to them.''

Orser, who coaches at the Toronto Cricket, Skating and Curling Club in Toronto, has witnessed that
celebrity firsthand. He and his former coach Uschi Keszler know how rabid fans can get.

Keszler said: ``Once, when we were somewhere in Asia, Brian blew his nose into a tissue during a
practice. A fan walked over and took it as a souvenir.''

Orser's experience with fans has kept him on guard for the fans who try to get too close. So, when
he received the multiple e-mail messages recently from Kim's teenage admirer, he contacted the
rink's security.

Receptionists were told to call security and maybe even the police if the teenager showed up, Orser
said. The fan's photograph, once sent to win Kim's heart, was posted at the rink's entrance like a
mug shot.

``I'm sure the guy is harmless,'' Orser said, ``but we didn't want to take any chances.''

\section{Stitching the Narrative of a Revolution}

\lettrine{I}{t}\mycalendar{Jan.'10}{26} was the height of the Cultural Revolution, but in the heart
of China's capital, in range of the prying eyes of foreign embassies, young Beijingers had embraced
the tenets of capitalism.

\myphoto{Stitching.the.Narrative.of.a.Revolution.jpg}{asants recited quotations from Mao's ``Little
  Red Book'' before toiling in the fields in a village near Beijing in July of 1967.}

Corrupted by dreams of profit, crowds of 500 or more were gathering every Sunday on a street in the
city's embassy district to ply a shameful trade. ``They are learning how to do business and raise
money,'' one city official wrote darkly. ``This is seriously harmful to the healthy growth of the
successors of the proletariat revolution.''

Such was the state of affairs in 1966, when selling pigeons at an impromptu street market was seen
as an obstacle to the triumph of socialism -- and, the official added, as a waste of bird feed, too.

The records on the Beijing pigeon market, like thousands of other Cultural Revolution documents, lay
silent for decades, deemed state secrets by a government hardly eager to highlight Mao's excesses.
But last year, China quietly opened the archives of selected declassified government files from that
era, in Beijing, Shanghai and Xi'an.

And so a veil has begun to lift on this and other prosaic stories of the Cultural Revolution -- some
sad, some funny, most humdrum to an extreme.

The files of the Cultural Revolution, which raged from 1966 until Mao's death in 1976, make up a
mere 16 of the 21,568 volumes that the Beijing Municipal Archives has made public in four separate
releases -- in 1996, 1997, 2001 and 2009. (The other files cover periods of Chinese history from
1906.) Stored in thick binders on library-style stacks, they can be viewed in the Municipal Archives
building, a spacious, modern structure with overstuffed chairs and a scholarly atmosphere on the
south side of the city.

The yellowing files give scant insight into those days' atrocities: the denunciations of parents by
children; the humiliation of intellectuals; the millions of lives ruined by Red Guards ordered to
remake society through upheaval. Mao's personality cult made him a living god, and armed violence
broke out over his affections. Everything was politicized. Many committed suicide.

Today, that era has been all but obliterated from the official history of the People's Republic, its
horrors glossed over in history books. While many younger Chinese know that the country passed
through a period of turmoil, scholars say, few have any idea of its wild extremes. Events that were
``earth shattering have now turned into words with vague and sketchy meanings,'' Chen Xiaojing, a
Communist Party official from the time, wrote in a carefully hedged account of his experiences, ``My
Cultural Revolution Years.''

Why the government is releasing some documents from the era is unclear. Archive officials declined
repeated requests for interviews. Experts say the files contain little if any material that
government censors would regard as incendiary.

``For people like me who have been studying the Cultural Revolution as a profession, it's better
than having nothing at all,'' said Xu Youyu, a historian and former researcher at the Chinese
Academy of Social Sciences. ``But the things I want to know are, for example, how many homes the Red
Guards had gone to raid and what they took out of each home. There's not a chance of finding those
things in these documents.

``If you air these things out, people may start asking why it happened. And this is not a question
that is directed only at 1966, but may be turned around and asked about the current situation in
China.''

Yet a picture of Chinese life 40 and 50 years ago does emerge from the archives. The files, some
nearly transparent and thin as one-ply tissue paper, include handwritten drafts of speeches, lists
of production quotas, song lyrics, government regulations and minutes of groups that studied Mao's
words. The texts embrace the political rhetoric of the day, in which all problems were succinctly
rendered into rhyming epithets.

The files apparently have been filtered for anything dealing with deaths and imprisonment, and they
describe a country still fervently Communist, and unrecognizable today. They narrate the story of a
country in the throes of madness, when ``Mao Zedong thought'' cured everything from truancy to
traffic jams to agricultural chemistry to illegal pigeon sales.

Consider: records from 1972, taken at a grade school outside Beijing, show that math students were
made to sing two revolutionary songs and study and discuss six Mao quotations for 25 minutes of each
class. The remaining few minutes were spent doing math.

In 1967, a report urged forming special groups at the provincial and city levels to ``use every
conceivable means to guarantee production'' each year of 13,000 tons of specially formulated red
plastic -- required for the covers of Mao's ``Little Red Book'' of quotations.

``The Conference on the Situation of the Special Plastic Used by the Works of Chairman Mao''
proclaimed that producing the plastic was ``our glorious political responsibility.'' To hold
everyone accountable, the conference produced a chart with a month-by-month breakdown of production
levels.

At times, the files veer perilously close to black, or perhaps red, comedy. In 1970, the annual
Representative Conference of the Enthusiasts of Chairman Mao's Works from the City Transportation
Bureau studied rush-hour bottlenecks created because workers were required to arrive early to study
Chairman Mao's works. The bottlenecks, the workers concluded, were the work of ``conservative
rightists and selfish departmentalism and other mistaken ideas.''

Yet there are also oblique hints of more sinister processes at work.

Many reports began with anecdotes of selfless revolutionary fervor. In one of them, Liu Chunnong, a
transportation security guard, recounted in 1968 how his dozen pet goldfish had been his pride and
joy. After a party meeting, he said, he took the fish outside and buried them alive. Raising
goldfish, he wrote, had been criticized as a petit bourgeois practice.

In a handwritten series of 1972 speeches, many of them heavily edited in pen, a teacher from
Beijing's outskirts recalled how his comrades ``patiently and delicately'' sought to reform a
teacher who was not a worker, but a member of the wealthy class. Rounds of criticism had little
effect, so the group chose to help him realize his mistakes through physical labor, by weeding
farmland.

``He pulled grass,'' the speech read. ``At first, he was squatting, but he couldn't handle it after
two days. Then he pulled the grass while kneeling. Finally, he did it while crawling.''

Party censors excised the tale of the exhausted teacher from the final draft of the speech.

\section{China to Move Against Local Lobbyists}

\lettrine{A}{ddressing}\mycalendar{Jan.'10}{26} a facet of political life all too familiar to
Americans, China's government is reported to have ordered the closure of thousands of ``regional
liaison offices'' -- in essence, lobbying firms -- that local governments and companies operate in
Beijing to curry favor with high officials.

Echoing another facet of political life all too familiar to Americans, not everyone is convinced
that the effort will succeed.

The government-affiliated magazine Outlook Weekly reported that Beijing officials have vowed to
shutter the liaison offices within six months in an effort to stanch what some analysts call a
culture of unalloyed corruption surrounding the institutions.

Liaison offices recently were featured in less-than-flattering light after it was reported that two
local governments from Henan Province last February spent more than \$96,000 to buy 777 bottles of
expensive Chinese liquor to entertain high officials.

That is but the latest in a string of scandals surrounding the offices, which have produced a number
of indictments in recent years on charges includingbribery and embezzlement. Chinese wags often
refer to the liaison offices as ``pao bu, xing jin,'' which ostensibly means ``run forward,'' but
also can be translated as ``go to the ministry and give money.''

By some calculations, Beijing boasts more than 5,500 local-government provincial offices, and
perhaps 5,000 more offices representing state-run corporations, associations and other entities.

The provincial-level offices alone employ close to 8,000 workers, housed in impressive bureaus
designed to convey the power and prestige of the areas they represent.

The central government directive would close the offices of state-run firms and governments at or
below the county level. The survivors, Outlook said, would be strictly regulated.

But the central government could be in for a fight. ``This is serious,'' Russell Leigh Moses, a
Beijing-based analyst of central government affairs, said in a telephone interview on Monday.
``These offices are strong symbols of provincial sovereignty, and by seeking to shut them down, the
central government is trying to keep regional officials from lobbying too hard for local
interests.''

China's leaders have used increasingly dire language in recent years to inveigh against government
corruption, which they liken to a disease that is eroding the strength of Communist governance.
Local government corruption is a common subject of attack by local journalists and internet surfers.
But with some exceptions, higher-level corruption has largely been excluded from government
campaigns.

The South China Morning Post, a Hong-Kong based newspaper, stated Monday that Beijing is awash in
liaison offices because state power is so concentrated that virtually anyone seeking influence needs
a representative in the national capital.

``In a democratic country, there is a fairly transparent regime in place to decide where a
government project goes and how much expenditure it will receive, even with some lobbying,'' Ren
Jianming , deputy director of Tsinghua University's Anticorruption and Governance Research Center,
was quoted as saying. ``But in China, it often comes down to networking and even some shady ways to
court higher authorities. So the liaison offices are indispensable in this sense.''

But the success of the government directive is anything but guaranteed. The Outlook article noted
that China's central government has been talking since at least 2006 about reining in the offices,
with little success.

One problem is that the offices represent China's provinces, which are also represented on the
nation's principal governing body, the central committee of the Chinese Communist Party. Provincial
officials at the national level are hardly enthusiastic about curtailing the power of officials in
their home regions.

\section{With Apple Tablet, Print Media Hope for a Payday}

\lettrine{W}{ith}\mycalendar{Jan.'10}{26} the widely anticipated introduction of a tablet computer
at an event here on Wednesday morning, Apple may be giving the media industry a kind of time machine
-- a chance to undo mistakes of the past.

Almost all media companies have run aground in the Internet Age as they gave away their print and
video content on the Web and watched paying customers drift away as a result.

People who have seen the tablet say Apple will market it not just as a way to read news, books and
other material, but also a way for companies to charge for all that content. By marrying its
famously slick software and slender designs with the iTunes payment system, Apple could help create
a way for media companies to alter the economics and consumer attitudes of the digital era.

\myphoto{With.Apple.Tablet.Print.Media.Hope.for.a.Payday.jpg}{An Apple Store in Chicago. Apple's
  anticipated tablet is expected to help media companies charge for their content}

This opportunity, however, comes with a sizable catch: Steven P.~Jobs.

Mr.~Jobs, the chief executive, made Apple the most important distributor of music by imposing its
own will on the music labels, bullying them into accepting Apple's pricing and other terms. Apple
sold lots of music, but the music labels claimed that iTunes had destroyed the concept of the album
and damaged their already deteriorating bottom lines.

With the new tablet, media companies could be submitting themselves to similar pricing restrictions
and sacrificing their direct relationship with customers to Apple.

For now, at least, the technology and media industries are looking at the brighter side. ``Steve
believes in old media companies and wants them to do well,'' said a person who has seen the device
and is familiar with Apple's marketing plan for it, but who did not want to be named because talking
about it might alienate him from the company. ``He believes democracy is hinged on a free press and
that depends on there being a professional press.''

Part of the media industry's high hope for the tablet comes from descriptions of the device from
analysts and others who have been briefed on it.

It will run all the applications of the iPhone and iPod Touch, have a persistent wireless connection
over 3G cellphone networks and Wi-Fi, and will be built with a 10-inch color display, allowing
newspapers, magazines and book publishers to deliver their products with an eye to the design that
had grabbed readers in print.

Their optimism for the tablet also stems from consumers' willingness to spend money using mobile
devices. In the last decade, while people downloaded music illegally to their desktop computers,
they happily paid small amounts of money on their cellphones to download ring tones and send text
messages.

The iPhone has provided further proof that the economics of mobile devices are unique: the Apple App
Store is expected to generate an estimated \$1.4 billion this year, according to an analysis by
Piper Jaffray.

``The iPhone was a harbinger,'' said Trip Hawkins, a founder of Electronic Arts and now chief
executive of Digital Chocolate, which makes games for cellphones. ``When you have a device that is
this convenient and fun for consumers to use, you can get a lot more people interested in paying for
and engaging with the content. Big media companies should be all over this like a cheap suit.''

Indeed, they already are. The New York Times Company, for example, is developing a version of its
newspaper for the tablet, according to a person briefed on the effort, although executives declined
to say what sort of deal had been struck.

On Monday, The Times also announced that its media group division had created a new segment for
``reader applications,'' and named Yasmin Namini, the senior vice president for marketing and
circulation, to head it. Executives said the timing was coincidental, prompted not by the Apple
device specifically, but by the growing importance to The Times of electronic reading devices in
general.

Two magazine publishers, Cond\'e Nast and Time Inc., have also created mockups of their magazines
for tablets, even before such devices have hit the market. ``Apple upended the smartphone market
with the introduction of the iPhone, and it's likely that they will, if they enter the tablet
market, lead the pace there,'' said Thomas J.~Wallace, editorial director of Cond\'e Nast. He said
that ``2010 is going to be the year of the tablet, and we feel we are in a very good position for
it.''

To successfully sell their material on the coming wave of tablets from Apple and other hardware
makers like Hewlett-Packard, media companies may first have to adjust other parts of their digital
strategies -- so consumers don't simply use the tablet's browser to get the same content free on the
Web.

Such shifts are under way.

In October, The Wall Street Journal, which is owned by the News Corporation, began charging for
access for certain elements of its iPhone application. Esquire and GQ have taken steps toward
charging for digital content, offering iPhone versions of their magazines for \$2.99 for each issue.

The December issue of GQ was downloaded from the app store almost 7,000 times, and twice as many
times for its January issue. Last week, The New York Times announced plans to begin charging, by
next year, frequent Web site visitors who are not also newspaper subscribers to read the online
version.

Media companies may have to swallow hard before tethering their futures to any high-tech company,
let alone Apple. Many publishers believe their economic health depends on finding a direct line to
their customers, and it is not clear whether Apple -- and other aggregators of Internet content --
will allow that.

Magazine publishers, for example, maintain sophisticated databases about their customers, which lets
them cross-sell products, renew subscriptions and entice advertisers with statistics about their
wealthy readers. A big part of the business is automatic renewals charged to credit cards.

But when magazine publishers sell applications through the iTunes store, they do not get credit card
information or even the name of the buyer.

However, Apple, which makes most of its money selling devices, not content, has shown itself in some
cases to be a more benevolent warden of online content, than, say, Amazon.com. Unlike Amazon with
the Kindle, Apple allows application makers to set their own prices; some, like The Financial Times,
give away applications for the iPhone, but then bill customers directly for repeat use.

Nevertheless, concern over preserving the customer relationship is one reason, late last year, that
major publishers including Time, Cond\'e Nast, Meredith, the News Corporation and Hearst announced
they had formed a consortium, called Next Issue Media, that plans to run its own online store
selling digital issues and collecting consumer information.

``It's fundamental to the business model of publishers,'' John Squires, the interim managing
director of the consortium, said last month. ``We've always enjoyed an opportunity to know exactly
where our consumers are, and be able to market other products to them. It's a very key issue for the
founding members of this business.''

One branch of big media whose fortunes may not be lifted by an Apple tablet, at least initially, is
the TV business. Apple has also talked to television networks about offering access, for a monthly
fee, to a selection of their hit shows, bypassing traditional distributors.

But perhaps smarting from their experiences with Apple, many of the old-line media companies -- NBC
Universal, Viacom and Discovery among them -- shrugged at (or totally dismissed) Apple's plans for a
TV subscription package, according to executives briefed on the talks. A person briefed on Apple's
plans confirmed that such a subscription video option was not part of any immediate offering.

\section{As China Rises, Conflict With West Rises Too}

\lettrine{A}{s}\mycalendar{Jan.'10}{27}recently as 2008, when China was still an emerging economy
eager to put its best foot forward for Western consumers, it lifted censorship, at least
temporarily, on several Web sites before the Beijing Olympics. At the same time, it responded to
pleas from U.S.~and European politicians to cooperate on several other fronts.

These days, China is no longer emerging. It has emerged -- sooner and more assertively than had been
expected before the wrenching global financial crisis, which badly damaged all the established
industrial powers, from the United States to Europe and Japan.

Its currency, the renminbi, is frozen at an undervalued level, and Internet controls are stricter
than ever -- even as Google, one of America's most prominent companies, threatens to leave.

\myphoto{As.China.Rises.Conflict.With.West.Rises.Too.jpg}{A worker makes preparations in the
  Congress Center one day before the opening of the World Economic Forum in Davos, Switzerlan}

The severe recession has fast-forwarded history, catapulting an unprepared world into a period of
uneasy cohabitation between the United States, the No.~1 economy, and its eventual successor.

``China is the West's greatest hope and greatest fear,'' said Kristin Forbes, a former member of the
White House Council of Economic Advisers and one of hundreds of top officials and executives
flocking to this winter resort for the annual World Economic Forum, which is taking place Wednesday
through Sunday.

``No one was quite ready for how fast China has emerged,'' said Ms.~Forbes, a professor at the
Massachusetts Institute of Technology. ``Now everyone is trying to understand what sort of China
they will be dealing with.''

For the first time, economists point to Chinese spending -- not the U.S.~consumer -- as the key to a
global recovery. China's gross domestic product could overtake that of the United States within a
decade, one report predicted this month, while others speculated about when the renminbi might start
to challenge the dollar as the world's reserve currency.

And as developing countries everywhere look for a recipe for faster growth and greater stability
than that offered by the now-tattered ``Washington consensus'' of open markets, floating currencies
and free elections, there is growing talk about a ``Beijing consensus.''

China's rise will be on prominent display in Davos this week, with the biggest Chinese delegation in
the World Economic Forum's history. The local Chinese restaurant has been fully booked since early
January. The 54 Chinese officials and executives -- including the presidents of the country's
sovereign wealth fund and export-import bank -- were expected not only to rub shoulders here but
also, as one put it bluntly, to ``go shopping.''

When the United States was snapping at the heels of the British empire, the global hegemon of the
early 20th century, the situation caused plenty of friction, even though both countries spoke the
same language, shared similar cultures and were liberal democracies.

China, in contrast, is a Confucian- Communist-capitalist hybrid under the umbrella of a one-party
state that has so far resisted giving greater political freedom to a growing middle class. Now its
ascendancy is about to set off what many officials and experts see as a backlash on both sides of
the Pacific.

``It's not surprising that China's remarkable economic rise would be unsettling to some,'' said
Pascal Lamy, the director general of the World Trade Organization.

So far, the backlash against China has been largely rhetorical. Stephen Roach, the Asia chairman of
Morgan Stanley, counts 45 anti-China legislative measures introduced in the U.S.~Congress between
2005 and 2007. None passed.

That could change, as tricky midterm elections loom in the United States and politicians there and
in Europe become more outspoken in blaming China's currency peg to the dollar, which gives its
industries a competitive edge, for rising joblessness at home.

Some targeted tariffs have been imposed in recent months. Washington has penalized imports of
Chinese tires and coated paper products. Both the United States and the European Union are
restricting Chinese steel.

But none of those measures go as far as climate change proposals in France and the United States,
which call for border taxes on products from countries -- China in particular -- that do not accept
higher costs for carbon emissions in producing energy and making goods. If ``the U.S.~opts for
friction,'' Mr.~Roach said, ``the Chinese can be expected to respond in kind.''

China has its own version of political jockeying. Several foreign companies already complain that
doing business in China has become more difficult. Lured until a few years ago by tax rates less
than half of those applying to Chinese companies, executives now cite an increase in red tape and a
growing number of ``buy China'' mandates from government procurement offices.

The standoff with Google has illustrated the difficulties foreign business faces in China. It has
also starkly raised the question of who will have the upper hand in future negotiations.

``The operating environment is tougher than ever for Western companies,'' said James McGregor, head
of the government relations committee of the American Chamber of Commerce in China. ``But unlike
Google, most Western companies also need China more than ever.''

China is the biggest recipient of foreign direct investment in the world: 450 of the Fortune 500
companies have business presences there, and many of those still reeling at home are doing brisk
business in China. ``G.M. is hurting anywhere else, but here things are quite profitable,''
Mr.~McGregor said.

Business interests in China could make it harder for Western politicians to lash out. ``It's a
situation the U.S.~was in for a long time,'' said Ms.~Forbes, the M.I.T. professor. ``Many people
didn't like U.S.~policy, but you had to be in the U.S.~market.''

If business executives are looking to China for its low manufacturing costs and sizable market,
political leaders are studying a state perceived to have found a recipe for lifting millions out of
poverty with fast growth, even if that means a stiff measure of domestic repression. ``You hear more
and more people talking about a Beijing consensus,'' Ms.~Forbes said.

But what exactly is the Beijing consensus? Some see it as a form of economic management with greater
government involvement that is on the rise across the world. Others interpret it to mean more
strictly controlled capital markets, which have made a re-appearance even in previously open
countries like Brazil. Policy makers in Malaysia and Dubai focus on replicating China's special
economic zones, which afford generous terms to foreign investors in manageable geographic areas.

Some suggest that China's lack of democracy is an advantage in making unpopular but necessary
changes. ``It is more challenging for democratic systems because every day they come under public
pressure and every short period they have to go back to the polls,'' said Victor Chu, chairman of
First Eastern Investment Group in Hong Kong, the largest direct investment firm in China. ``China is
lucky to have the ability to make long-term strategic decisions and then execute them clinically.''

With China's rising clout, the West has less leverage over Beijing. When China was seeking to join
the World Trade Organization a decade ago, it accepted compromises to U.S.~and European demands. At
climate talks last month in Copenhagen, however, China blocked a comprehensive deal and refused to
go beyond its earlier promises. Portrayed as a deal breaker in the Western media, at home it was
celebrated as the country that stopped the West from imposing its terms on developing countries,
Mr.~Chu said.

Western diplomats complain about the way Beijing is dragging its feet more than Moscow on sanctions
on Iran's nuclear program and is propping up unsavory regimes across the world in its hunt for the
natural resources to power its growth.

Some say Chinese officials are using their country's \$2.4 trillion in foreign currency reserves as
a bargaining chip, knowing that any hint of reducing those reserves would rattle currency markets.

``As China is emerging on the global stage with unprecedented power and influence,'' said David
Shambaugh, a professor of political science and international affairs at George Washington
University who is in China as a Fulbright scholar, ``it is not proving to be the global partner the
United States and E.U. seek.''

In the world of power politics, that is not particularly surprising. Like many Western countries,
China will act only when it is in its interest.

Mr.~Chu of First Eastern Investment said he expected China to resume a gradual appreciation of the
renminbi later this year, not because Washington was lobbying for it but because signs of
inflationary pressure and bubbles in the Chinese credit and housing markets were mounting. This
month, the Chinese authorities raised interest rates and moved to curtail bank loans.

Kenneth Rogoff, an economics professor at Harvard University who just spent two weeks in China,
warns that the country will face its share of economic troubles in the years ahead. But that will
not change the underlying trend, he said.

While China remains much poorer than the advanced industrial powers of the West on a per-capita
basis, its rapid growth should enable it to pass Japan this year as the world's second-largest
economy.

A new report by PriceWaterhouseCoopers predicts that China could overtake the United States as the
largest economy as early as 2020. In 2003, Goldman Sachs made waves by suggesting that the Chinese
G.D.P. might match that of the United States by 2041. Five years later, the forecast was revised to
2027.

According to Mr.~Rogoff, over the next four decades or so, the Chinese renminbi will gradually come
to rival the dollar as the world's leading reserve currency, making China's response to its
increasingly central role in the global economy critical.

The risk, Mr.~Shambaugh of George Washington University said, is that ``the world will be asking
more and more of China but getting less and less in return.''

\section{Dalai Lama's Envoys to Resume Talks With China}

\lettrine{E}{nvoys}\mycalendar{Jan.'10}{27} of the Dalai Lama, the spiritual leader of the Tibetans,
arrived in China on Tuesday to resume a long-running series of negotiations over Tibet that have so
far been fruitless.

The current set of negotiations will be the ninth since 2002 between the Dalai Lama, whom China
accuses of being a dangerous ``splittist,'' and the Chinese government, which has flooded large
areas of Tibetan regions in western China with security forces since a widespread uprising there in
March 2008.

The Dalai Lama, 74, who lives in exile in the Indian hill town of Dharamsala, says he does not want
independence for Tibetan regions, but rather greater autonomy -- the Tibetans should be able to make
their own policy regarding religious practice, education and immigration to the regions, he says.
Tibetans are anxious over the large numbers of ethnic Han, who dominate most of China, moving to the
Tibetan plateau to seek their fortunes.

The series of negotiations between the Dalai Lama and China broke down in November 2008, after
Chinese officials strongly rejected a proposal for genuine autonomy presented to them. During the
eighth round of talks, the two envoys who handed over that proposal were Lodi G.~Gyari, from
Washington, and Kelsang Gyaltsen, of Switzerland. Both are representing the Dalai Lama again in this
round.

The envoys will return to India in early February, the Tibetan government-in-exile said in a
statement on its Web site.

The United Front Work Department, the part of the Chinese Communist Party that officially deals with
ethnic policy, said in a written statement on Tuesday: ``We hope the Dalai Lama side will cherish
this opportunity and eagerly respond to the demands by the central government.''

Tibet remains one of the most delicate issues between China and the United States. President Obama
has said he will meet with the Dalai Lama early this year. Last fall, Mr.~Obama declined to meet
with the Dalai Lama while the Dalai Lama was visiting the United States. The move was widely seen as
an attempt by Mr.~Obama to improve United States-China relations, but Chinese leaders continued to
press the president on Tibet during his first state visit to China in November. State Department
officials quickly abandoned a new phrase that had been coined to describe the conciliatory approach
-- ``strategic reassurance'' -- and which had been criticized by some Americans as being tantamount
to appeasement.

Late Monday, the State Department said it welcomed the latest round of negotiations between the
Dalai Lama and China.

``The United States strongly supports dialogue between China and the Dalai Lama's representatives to
address longstanding differences,'' said Philip Crowley, a State Department spokesman, in a written
statement. ``The administration hopes this meeting will produce positive results and provide a
foundation for future discussions to resolve outstanding issues.''

Though Tibetans uniformly revere the Dalai Lama, there are many, especially younger ones, who
criticize the Dalai Lama's approach to China. They say that the Tibetan exile movement should be
seeking independence for Tibet, not accommodation under Chinese rule. The Chinese government has no
intention of giving ground, they say, and is just using the series of negotiations to stall for time
until the Dalai Lama dies.

``Until we see visible change on the ground in Tibet and some real give-and-take on the part of the
Chinese government, we can be confident that the resumption of the dialogue is nothing more than a
delaying tactic designed to mute international criticism, especially in the lead up to the expected
meeting between Obama and the Dalai Lama,'' said Tenzin Dorjee, executive director of Students for a
Free Tibet, a group based in New York that advocates for Tibetan independence.

The Dalai Lama's spokesman, Tenzin Taklha, said Monday that the talks were part of an ``important
process of trying to find a mutually agreed solution,'' according to Agence France-Presse.

Last week, China's top leaders met to draw up plans for governing Tibet. State news organizations
said the leaders determined that continued economic development would bring ``stability'' to the
region.

\section{Obama Vows Action on Ailing Economy}

\lettrine{W}{ith}\mycalendar{Jan.'10}{28} his legislative agenda stalled and unemployment stuck in
double-digits, President Obama plans to use his first formal State of the Union address Wednesday
night to outline his efforts to rebuild the nation's economy and put more Americans to work.

\myphoto{Obama.Vows.Action.on.Ailing.Economy.jpg}{President Obama walked to the Oval Office on
  Wednesday afternoon.}

As part of that revitalization, aides said he would vow to fight to salvage the comprehensive
overhaul of the health care system that was his domestic priority until it collapsed in Congress. He
will also use the speech to make some pointed attacks on Washington lobbyists, according to early
excerpts of his prepared remarks.

In those remarks, Mr.~Obama paints a grim picture of the nation's economic challenges and the dire
state of health care costs.

``By the time I'm finished speaking tonight, more Americans will have lost their health insurance,''
he says in the excerpts. ``Millions will lose it this year. Our deficit will grow. Premiums will go
up. Co-pays will go up. Patients will be denied the care they need. Small business owners will
continue to drop coverage altogether.''

But, he adds: ``I will not walk away from these Americans. And neither should the people in this
chamber.''

Mr.~Obama's aides have said he will concede to having made mistakes during his first year in office
and call on his Republican critics to redirect their energy and help him rebuild the economy and
change the tone and culture in Washington.

``Rather than fight the same tired battles that have dominated Washington for decades, it's time for
something new,'' the president says in the excerpts. ``Let's try common sense. Let's invest in our
people without leaving them a mountain of debt. Let's meet our responsibility to the people who sent
us here.``

The economic and political climate for this set piece of Washington theater is far more abrasive
than Mr.~Obama might have imagined when he was elected the 44th president in 2008 and declared that
``change has come to America.''

The nation has been mired in the worst economic downturn since the Great Depression in the 1930s,
prompting him to focus on how he will create new jobs and respond to the mounting insecurities and
frustrations of working families and the middle class.

His moves will include some tax credits, subsidies for expenses like day care and student loans and
a three-year spending freeze for discretionary spending. He also plans to appoint a bipartisan
commission to make recommendations to reduce the nation's long-term debt.

But he cautions: ``We have to recognize that we face more than a deficit of dollars right now. We
face a deficit of trust – deep and corrosive doubts about how Washington works that have been
growing for years.''

He adds: `` To close that credibility gap we must take action on both ends of Pennsylvania Avenue to
end the outsized influence of lobbyists; to do our work openly; and to give our people the
government they deserve. That's what I came to Washington to do.''

Mr.~Obama also plans to ask Congress to repeal the ``don't ask don't tell'' law that prevents gays
from serving openly in the military.

Presidential administrations often use early excerpts of a major speech like this to try to shape
the real-time analysis online and on television and in early editions of newspapers. Mr.~Obama's
first excerpts focus on lobbyists and the corporate culture in Washington that he had vowed to
change, though it is possible they may not reflect the major thrust of the speech that he delivers
when he arrives in the House chamber.

Unemployment is now higher than when Mr.~Obama took office, and polls show that Americans think the
country is on the wrong track by a two-to-one margin. The president's top priority, a health care
overhaul, has been put on the back burner by his fellow Democrats, who are feeling rattled after the
recent Republican victory in the Senate race in Massachusetts.

Mr.~Obama himself has acknowledged that the health bill smacked of the very kind of backroom
deal-making he campaigned against, and in the days leading up to the address he indicated he
intended to use it to offer a mea culpa, if not for acceding to those deals, then for failing to
change the culture of Washington. White House officials said in advance of the speech that Mr.~Obama
would use it to talk about changing the way Washington does business.

But the big issue, for Mr.~Obama, is the economy, and the White House signaled he would devote a
major chunk of the speech to it.

``For a decade, American families and businesses have experienced an erosion of the economic
security that has been at the core of our prosperity,'' said a memorandum distributed by the White
House to Democrats on Capitol Hill Wednesday afternoon, describing the broad outlines of Mr.~Obama's
message. ``Tonight the president will lay out his vision on what steps are required to restore that
security and how Washington should do its work in a way that reflects the best of our people and our
democracy.''

The speech is Mr.~Obama's third to a joint session of Congress. In recent days the White House had
made the broad outlines of the speech clear: it would put a heavy emphasis on job creation, and
Mr.~Obama would calling on Congress to pass a jobs bill that would spur investments in small
businesses and in green jobs and clean energy.

He would outline a series of modest proposals intended to provide aid to the middle class, including
new tax credits for child care and a cap on student loan payments for recent graduates. He would
call for a three-year freeze on spending on many domestic programs, and announce the creation of a
bipartisan commission that would recommend tax increases and spending cuts with the aim of bringing
down the budget deficits and slowing the growth of the national debt.

Allies said he intended to put the three major elements of his agenda – jobs, healthcare and a new
energy policy – into broader context, to describe how they fit together to revive the sagging
economy.

``Rescue, Rebuild Restore – a New Foundation for Prosperity,'' read the headline of the White House
memorandum, summing up Mr.~Obama's message.

That phrase – ``new foundation'' – harks back to the earliest days of Mr.~Obama's presidency, when
the White House used it to try to provide a broad narrative for Mr.~Obama's policies. He used it in
his first address to a joint session of Congress last February, declaring, ``Now is the time to act
boldly and wisely, not only to revive this economy but to build a new foundation for lasting
prosperity.'' He used it again during a speech last April on the economy.

``He's got to get that narrative back,'' John D.~Podeseta, a close ally of Mr.~Obama's who also
worked as chief of staff to President Bill Clinton, said in an interview Wednesday. ``That's what he
lost over the last eight months.''

But while Mr.~Obama may be able to recapture his narrative, even Democrats concede he is unlikely to
recapture the sense of goodwill and bipartisan spirit that he brought to office with him one year
ago. Republicans said on Wednesday that while they welcomed Mr.~Obama's emphasis on fiscal
responsibility and his call for a three-year spending freeze on a portion of the domestic budget,
but they also came out criticizing the president and his agenda even before the speech was
delivered.

``We don't need another lofty speech from his historic well from the president of the United
States,'' Representative Mike Pence, the Indiana Republican and chair of the House Republican
caucus, told reporters. ``Mr.~President, set aside your big-government plans to take over health
care; send us a budget that reduces spending and reduces taxes on the American people. Mr.
President, deeds, not words.''

\section{Haitians in the U.S.~Find Obstacles in Immigration Process}

\lettrine{M}{arie}\mycalendar{Jan.'10}{28} Violande Guerrier-Cavalier arrived in Florida from Haiti
on Jan.~16, with little more than her feverish infant son, Marcley, his tiny legs in casts because
of a birth defect. She left her husband and four other children behind, living in the yard outside
their broken house.

Because Marcley is a United States citizen, born here, his mother was allowed to evacuate with him
after the earthquake, and can stay in the United States for six months. But Ms.~Guerrier-Cavalier
has few possessions, no money and no way of legally earning any. She can neither support herself and
Marcley, nor send money back home at a time when Haitians are desperate for help from outside. And
it is unclear where she will go when the six months are up.

``I came with the intention of being able to support my family,'' said Ms.~Guerrier-Cavalier, 45,
who for now is staying with a friend. ``To see my kids not being able to eat, that hurts really
badly.''

Federal officials said they were still trying to figure out what to do for parents like
Ms.~Guerrier-Cavalier. And they are struggling to deal with an array of other complexities involved
in helping both earthquake victims that have made it to the United States and those they left
behind.

About a week after the earthquake, the United States extended temporary protective status to
Haitians who were already in the United States on the day of the quake, allowing them to stay
legally and to work.

But hundreds of Haitians like Ms.~Guerrier-Cavalier, who were allowed to come here after the
earthquake as escorts for American citizens who are minors, are not eligible for that status.

They also are not eligible for a status known as humanitarian parole, which would allow them to
apply for permission to work, and to renew their status in one-year increments. Some Haitian orphans
were granted that status, allowing them to come here before their adoption paperwork was finished.

Federal officials said they had realized their mistake and, last weekend, began to give incoming
parents humanitarian parole instead of tourist visas. But at least one mother who arrived on Sunday
had only a tourist visa, said Gepsie Metellus, the director of Sant La Haitian Neighborhood Center
in Miami. And the change came too late for Ms.~Guerrier-Cavalier, and others like her. Humanitarian
parole cannot be given to someone already in the United States.

In the aftermath of the earthquake, temporary protected status was made available only to Haitians
already in the United States because officials were seeking to deter Haitians from trying to get to
the United States by boat. No such mass exodus has materialized, but other problems with the program
have surfaced.

One reason officials cited for extending temporary protected status to Haitians was so they could
get jobs and send more money home. But applicants are being asked to scrape together nearly \$500 in
government fees, money that could be going instead to family members in Haiti.

Legal advocates said many of the estimated 100,000 Haitians eligible for temporary protected status
have not been working and cannot afford the fees.

``The No.~1 problem is the money,'' said Randolph P.~McGrorty, the chief executive of Catholic
Charities Legal Services in Miami, which collected more than 2,000 applications in the last week,
with many still lacking the required money order.

Mr.~McGrorty said that it was possible to get the fees waived, but that such requests would delay
applications. ``So far, they have refused to grant a blanket fee waiver,'' he said.

Applicants pay \$50 for temporary protected status, \$80 for a background check and \$340 for
permission to work. Matthew Chandler, a spokesman for the Department of Homeland Security, said that
the fees were set by Congress and that only Congress could allow a blanket fee waiver.

Though federal officials have said they would be ``generous'' when reviewing individual applications
for fee waivers, it is not clear exactly what it takes to qualify.

``It depends on a case-by-case basis,'' said Chris Bentley, a spokesman for United States
Citizenship and Immigration Services, which runs the program. ``There's no one answer, that here's
what you have to submit to be able to get it.''

Maricame Aristide, 63, who has supported herself as a live-in maid since her application for
permanent residency was turned down, came to St. Clements Catholic Church in Broward County on
Monday for help with her application. She said that even if she had that kind of money, she would
rather send it to her orphaned grandson, now living on the streets of Port-au-Prince with
Ms.~Aristide's sister.

Evacuees continue to pour into the United States, many in need of basics like clothes and diapers.
As of Monday morning, 15,666 people had been evacuated from Haiti to the United States. Of those,
12,497 were American citizens, according to State Department figures. The rest were legal permanent
residents or the adult escorts of minor citizens like Ms.~Guerrier-Cavalier, or like Darline Chery,
35, and her husband, Wladimyr, who managed to get their three daughters, all born in the United
States, to Fort Lauderdale.

The Chery family, whose members have only tourist visas, is staying in the one-bedroom apartment of
a cousin. The oldest daughter, who is 6, has already been enrolled in school. In Haiti, Mr.~Chery
was a travel agent and Ms.~Chery worked for a non-governmental organization. Without the ability to
work now, they are at a loss.

``We are all extremely traumatized,'' Ms.~Chery said, holding her 4-month-old baby and speaking in a
very quiet Creole.

Instead of being able to help themselves and relatives at home, parents with short-term visas have
been forced to ask for help themselves. Ms.~Guerrier-Cavalier and the Cherys have already applied
for Medicaid benefits and food stamps for their children.

``My husband is doing everything that can be done,'' Ms.~Chery said. ``Because when you have a
family and you're the parent, you're supposed to be taking care of them.''

\section{Venus Is Out, but Serena Perseveres}

\lettrine{T}{here}\mycalendar{Jan.'10}{28} will be no Williams sisters showdown in the semifinals of
the Australian Open. Li Na made sure of that Wednesday by stunning sixth-seeded Venus Williams, 2-6,
7-6 (4), 7-5, in a grinding, error-filled marathon that should have been over after two sets. Li now
joins her compatriot Zheng Jie in the semifinals, the first time two Chinese players have advanced
that far in a Grand Slam.

\myphoto{Venus.Is.Out.but.Serena.Perseveres.jpg}{Li Na defeated sixth-seeded Venus Williams, 2-6,
  7-6 (4), 7-5 on Wednesday.}

Serena Williams, the defending champion here and the top seed, struggled as well before beating back
a stiff challenge from seventh-seeded Victoria Azarenka, 4-6, 7-6 (4), 6-2. Williams trailed by 4-0
in the second set and was in danger of being eliminated when she stormed back to force a tie
breaker.

``I didn't expect to win when I was down, 0-4,'' Williams said. ``I thought, oh well, at least I'm
still in doubles.''

Azarenka knew she made a mistake by becoming tentative and not putting Williams away.

``The only thing I did is I probably backed up a little bit -- I gave her more time to go for her
shots,'' she said. ``I had chances. I had opportunities. It wasn't many, but I didn't take it.''

Roger Federer, the world's top player, turned back the hot-playing Nikolay Davydenko, 2-6, 6-3, 6-0,
7-5. The sixth-seeded Davydenko arrived at Melbourne Park with a 14-match winning streak, which
included two victories over Federer. He was flawless in winning the first set, and had Federer down,
3-1, in the second. Federer then won 13 games in a row, cutting the heart out of Davydenko, who
sprayed balls wide and long and grew frustrated as he made 50 unforced errors.

``I know he goes through his phases, and some of them are rock solid,'' Federer said. ``I was in a
tough situation at 6-2, 3-1 down and 15-40 on my serve. I knew I wasn't looking very good, you know.
But that's the beauty of best-of-five sets.''

Federer will face No.~10 Jo-Wilfried Tsonga of France, who, in a rematch of the 2008 Australian Open
final, outlasted an ailing Novak Djokovic, 7-6 (8), 6-7 (5), 1-6, 6-3, 6-1, in the other men's
quarterfinal. Djokovic, who was seeded third, said that he had an upset stomach before the match.
``When you lose a lot of fluids, your engines stop working, and that's how I felt,'' he told The
Associated Press.

Venus Williams looked crisp while winning her first set and served for the match in the second. Li,
who was down, 5-3, refused to quit, breaking Williams and forcing a tie breaker that she won by
blasting powerful forehands and whistling crosscourt backhands.

``Venus played aggressive in the first set,'' Li said. ``She didn't miss a lot of balls. I was
feeling more pressure in the first set. Then second set I was feeling a little bit better, but still
was, like, 5-3 down. Then I just try to play more ball back. Maybe have chance and then I do it.''

The third set was about perseverance. Williams held her serve in the first game, but she and Li
broke each other the next six games. Williams struggled with her serve, hitting only 56 percent of
her first serves and committing 11 double-faults.

Li, seeded 16th, broke Williams once more in the 11th game, but she needed three match points before
she was able to join Zheng, who is unseeded, in the semifinals. It will be the first time two
Chinese players have advanced that far in a Grand Slam.

``I just play my tennis; I know I lose first set but thanks for the coach -- he never gave up,'' the
27-year-old Li said, beaming to the grandstands at her husband, Jiang Shan; her coach Thomas
Hogstedt; and the trainer Richard Sutton.

Jiang and Li met when Li was 16 and married in 2006.

``Sometimes he talk too much, but nice guy,'' she said of Hogstedt.

The Australian Open has been vexing Venus Williams lately. Last year she lost in the second round.
She reached the semifinals in 2001. Her best finish was in 2003, when she lost in the final to her
sister Serena.

She knew that she was a match point away from a rematch in the semifinals. ``I felt confident
throughout the match, even in the third,'' Williams said. ``In tennis you have to close it out. It's
not like there's a clock ticking and then suddenly it's over. You just have to close it out. I
didn't do that today.''

Williams committed 53 unforced errors, and every time she tossed the ball in the air for her serve,
it was an adventure. She won only a third of her 46 second serves.

``I think I just started rushing my shots a little too much instead of just taking my time and
really just hitting a clean ball,'' Williams said. ``Sometimes I was able to back off. But, you
know, sometimes it's so easy to get so eager.''

Williams looked worse for the wear in the third set. This was her first tournament in three months.

``You know, she hits a good ball,'' she said. ``She hits a low ball. Sometimes it's not as easy to
move the low ball around as much. You know, we both did a lot of moving in the third set.''

Li, meanwhile, shook off her mistakes -- 57 unforced errors -- and hung around to reach the
semifinals of a major for the first time. She needs to get through Serena Williams if she wants to
land in the final, but for a day, at least, she wanted to live in the moment.

``I mean, this is the best day of my whole life,'' she said, knowing full that she and Zheng had
made history. ``It's good for both players, it's good for Chinese tennis, and also good for the
fans.''

\section{As Hong Kong's Political System Stalls, So Does Its Democracy Movement}

\lettrine{T}{he}\mycalendar{Jan.'10}{28} political system in Hong Kong is increasingly paralyzed,
and street protests are growing more confrontational as public dissatisfaction on economic issues
and a lack of democracy is rising. At the same time, the pro-democracy movement here has splintered,
weakening its ability to press for changes.

\myphoto{As.Hong.Kong.s.Political.System.Stalls.So.Does.Its.Democracy.Movement.jpg}{Protesters
  scuffled with police officers during a protest outside Hong Kong's Legislative Council on Jan. 16,
  2010.}

Protesters, many of them young people proclaiming their interest in democracy, have opposed building
an expensive high-speed rail link to Shenzhen and Guangzhou in mainland China. They are also upset
that Hong Kong's mostly unelected legislators approved the measure.

The demonstrations also reflect frustration on the part of the pro-democratic parties in the former
British colony that accuse China of having delayed or backtracked on commitments it made in the
1990s to allow people to directly elect a majority of lawmakers in the territory's Legislative
Council.

Donald Tsang, the Beijing-backed chief executive of Hong Kong, has suffered a significant decline in
his approval ratings in recent polls in recent months as he has grappled with rising local
discontent.

Until recently, Hong Kong had a tradition of orderly political protests that were uncommonly polite
by international standards. When 500,000 people took to the streets in 2003 to oppose successfully
the introduction of stringent internal security regulations, the police did not make a single
arrest.

Protests in the last few weeks, however, have been coarser. Youths have shouted obscene curses at
police officers. Scuffles with officers have resulted in a series of arrests.

``We're just sick of going to rallies that political parties organize, and we hold our banners and
don't accomplish anything,'' said Christina Chan, a 22-year-old graduate student in philosophy who
was arrested at her home this month on suspicion of assaulting police officers at two rallies.
Released on bail of 500 Hong Kong dollars, or \$65, she has not yet been formally charged and has
denied any wrongdoing.

Under the terms of its transfer to Beijing's rule, Hong Kong retains broad civil liberties but also
a political system that gives much greater weight to the votes of the economic and social elite.
Analysts say there is limited opportunity for youths to vent their unhappiness with dwindling social
mobility, high unemployment, sharply rising university tuition and an urban planning process
dominated by real estate developers.

Young people have borne the brunt of competition from low-salaried employees in mainland China. They
face rising competition for jobs in Hong Kong itself as banks and other high-paying employers
increasingly hire mainland college graduates with family connections in Beijing.

``This has been building for months, and I think we're heading for even greater frustration,'' said
Michael DeGolyer, the director of the Hong Kong Transition Project, a coalition of academics who
have traced the territory's political evolution for 22 years.

But the five pro-democracy political parties in the political opposition in Hong Kong have split
deeply over tactics this winter, including a move by two of the parties this week to bring about
by-elections.

Five lawmakers from two pro-democracy parties submitted letters of resignation on Tuesday that are
to take effect at midnight on Thursday. Unless withdrawn before then, the resignations would prompt
by-elections that the five hope to turn into an informal referendum on introducing greater democracy
before the next elections in 2012, instead of waiting until 2017 or later, as Beijing officials have
demanded.

Audrey Eu, the leader of the pro-democracy Civic Party and one of the five resigning, said that she
believed that unhappiness with the city's economic troubles could be effectively channeled into
support for democratic reforms. ``People are beginning to see it is really tied to our political
system,'' she said.

But the other three parties -- including the Democratic Party, the largest in the pro-democracy
movement -- have questioned the wisdom of this and have chosen not to have any of their lawmakers
resign. Albert Ho, the chairman of the Democratic Party, said that he did not believe that
dissatisfaction with the economy would show up in by-elections.

If the resigning lawmakers win their seats again in the by-election, then the Beijing-backed
executive branch of the government here will dismiss the results as meaningless, Mr.~Ho warned. But
if any of the lawmakers loses his or her seat, then the government will seize upon the results as
evidence that public support for greater democracy is limited, he said.

``If we embark on this project, we would be in a no-win situation,'' Mr.~Ho said.

Mr.~Tsang, the chief executive, said in a statement on Tuesday night that the government would not
recognize the results of the by-elections as a referendum. He cautioned that ``Many see it as an
abuse of the by-election mechanism and a waste of public resources.''

Beijing's main representation here, the Hong Kong and Macao Affairs Office, condemned the
resignations as only likely to cause more ``social conflicts.''

Economic worries here have crystallized around a plan to spend 67 billion Hong Kong dollars, or
nearly 10,000 Hong Kong dollars per resident, to build a high-speed rail link across the border to
Guangzhou. Backers of the rail line, including the government and the city's economic elite, see the
link as essential to tying the city into the mainland's rapidly improving rail system and its
vibrant economy, growing at 10 percent a year.

Critics of the rail line call it expensive at a time when the government is effectively raising
university tuition, looking for ways to limit assistance to all but the poorest senior citizens and
still mulling how to introduce a minimum wage.

Half the Legislature is elected by the general public, but the other half is chosen by so-called
functional constituencies -- groups ranging from Hong Kong banks to the city's lawyers. Only 800
people in this city of seven million are allowed to vote for the chief executive, although most of
the 800 are in turn elected by the same functional constituencies that control half the Legislature.

Ms.~Eu said she believed that the by-elections could become an informal referendum on whether to
eliminate functional constituencies and introduce the principle of one person one vote.

\section{Toyota's Push for Fast Growth Put Quality in Question}

\lettrine{T}{oyota}\mycalendar{Jan.'10}{28} executives set an ambitious goal in 2002 to own 15
percent of the global auto industry by 2010, meaning it would surpass General Motors as the world's
largest carmaker. To get there, it would have to grow by 50 percent. It would have to build new
plants in the United States, China, and elsewhere in Asia, and introduce dozens of new models.

Toyota managed to win bragging rights as the world's biggest car company. But that focus on rapid
growth appears to have come at a cost to its reputation for quality, creating an opportunity for
others to potentially take back market share they lost to Toyota.

On Tuesday, after prodding by the Transportation Department, Toyota announced it was temporarily
stopping production and sales of eight models that make up more than half its annual sales in the
United States, after two major recalls since November, while it tries to fix a problem with their
accelerator pedals.

``There was always a question about how fast they could go,'' James P.~Womack, an author and expert
on Toyota's manufacturing methods, said of the automaker's growth. ``I'm sure they regret that they
stomped on the gas so hard.''

Mr.~Womack said Toyota was also paying for taking its eye off the message that has been central to
its marketing.

``When your whole deal was quality, every mistake is a big deal,'' he said.

Transportation Department officials in the United States said Wednesday that they had advised the
automaker to act quickly.

``The reason Toyota decided to do the recall and to stop manufacturing is that we asked them to,''
Raymond LaHood, the transportation secretary, said.

In an interview with WGN radio in Chicago, Mr.~Hood added, ``We were the ones that really met with
Toyota, our department, our safety folks, and told them, 'You've got to do the recall.' ''

Japanese media raised fears that Toyota's problems might damage the reputation of other Japanese
companies.

``The discrediting of Toyota could even destroy the world's trust in Japanese manufacturing, which
relies on its reputation for high quality,'' warned the Tokyo Shimbun, a daily newspaper.

But fears about lasting damage to Toyota may be overblown, given the short attention span of
consumers, said Jeffrey K.~Liker, a professor of engineering at the University of Michigan and the
author of the best-selling book, ``The Toyota Way.''

``This is an unfortunate series of events that aren't a good reflection of the health of the
company,'' he said.

``But if they aren't selling cars, nothing else matters.''

Toyota announced a similar halt to sales and production in Canada, where it has two assembly plants
in Ontario. It is considering the same steps in Europe, although it has not yet come up with a plan.

Halting production and sales of the eight models in North America represents the latest setback for
Toyota's president, Akio Toyoda, who has been buffeted by problems since taking the job seven months
ago.

He has already apologized for the company's losses, its bland cars and its overconfidence. Last
fall, Mr.~Toyoda said the automaker was ``grasping for salvation.''

At the time, it seemed an overstatement. Despite its red ink, Toyota still has ample cash, and it
gained market share in 2009 in the United States, where it ranks as the second-biggest player behind
G.M.

On Monday, before Toyota announced the sales and production shutdowns, the company said it expected
its global sales to grow by 6 percent in 2010, to about 8.27 million vehicles.

But as it has gained sales, Toyota has moved away from some of the business practices it adopted in
its years of slow but steady growth.

One example is its decision to buy parts from companies around the world, rather than from a small
group of Japanese suppliers that have been longtime partners. For example, the pedals in the
vehicles affected by the production and sales stoppages come from a supplier's Canadian plant.

The move to expand its supplier network was necessary to save money, given cost pressures on Toyota
and other manufacturers, said Ulrike Schaede, professor of Japanese business at the University of
California, San Diego. But the shift also makes it harder for Toyota to control quality.

``At one level, it's the right thing to do,'' Professor Schaede said. ``At another, what you pay is
what you get.''

In Japan, ``many of us weren't surprised over the big recalls,'' said Masahiro Fukuda, manager of
research at Fourin, a global automotive research company based in Nagoya, Japan. ``We were more
surprised that it took Toyota so long.''

For Toyota, the timing of this problem is working somewhat in its favor. The production and sales
halts are coming at one of the slowest times of the year, allowing the company to begin repairs
before the car market picks up in the spring.

On Wednesday, Mr.~LaHood said engineers at the Transportation Department were working closely with
the company to find solutions.

Mr.~Womack said an open question is whether consumers will be patient with the company, because they
could choose from a variety of well-built cars and trucks from many other companies.

``They're trying to regain their image of quality at a time when everyone has gotten so much
better,'' he said. ``If everyone else is building almost-perfect vehicles, you have to make
absolutely perfect products.'' Mr.~Womack said he was still rooting for Toyota to fix its problems.

``I'm saddened to see a company known for building good things make some unnecessary mistakes,'' he
said.

\section{Rule of Thumbs: Koreans Reign in Texting World}

\lettrine{H}{a}\mycalendar{Jan.'10}{28} Mok-min is feeling like a gunslinger these days. At the
English-language cram school she attends during the winter break, students jealous of her
international bragging rights line up to duel with her.

\myphoto{Rule.of.Thumbs.Koreans.Reign.in.Texting.World.jpg}{Ha Mok-min, right, with her mother, Kim
  Young-sook, in Seoul, South Korea. Ms. Ha and Bae Yeong-ho won a texting contest.}

``They come with their cellphones boasting they can beat me,'' said Ms.~Ha, 16, her deadpan manner
lending her the air of a champion accustomed to -- even weary of -- fame. ``I let them try.''

She and another young South Korean, Bae Yeong-ho, recently conquered the world with their thumbs.
Their Team Korea won an international competition held in New York this month to determine who can
send text messages the fastest -- and most accurately -- on a cellphone.

``When others watch me texting, they think I'm not that fast and they can do better,'' said Mr.~Bae,
17, a high school dropout who dyes his hair a light chestnut color and is studying to be an opera
singer. ``So far, I've never lost a match.''

In the New York competition he typed six characters a second. ``If I can think faster I can type
faster,'' he said.

The inaugural Mobile World Cup, hosted by the South Korean cellphone maker LG Electronics, brought
together two-person teams from 13 countries who had clinched their national titles by beating a
total of six million contestants. Marching behind their national flags, they gathered in New York on
Jan.~14 for what was billed as an international clash of dexterous digits.

To ensure a level playing field, LG handed out identical mobile phones -- one with a numeric keypad
and the other with a keyboardlike QWERTY pad -- weeks in advance for practice. The basic rule of the
competition: copy phrases streaming across a monitor correctly, with the required capitalization and
punctuation, as quickly as possible. Whichever language players chose, words were selected so that
each would type the same number of characters.

Then they went thumb to lightning thumb, in five battles with names like ``The Monsters' Swamp'' and
``Race of Death.'' When it was over, Ms.~Ha and Mr.~Bae were proclaimed the world's fastest texters.
An American team came in second, an Argentine team was third.

Since their return home with \$50,000 in prize money each, Ms.~Ha and Mr.~Bae have become something
approaching heroes to what Koreans call the ``thumb tribe'' -- youngsters who feel more comfortable
texting than talking.

Until his recent immersion in music studies, Mr.~Bae texted 200 to 300 messages a day. Ms.~Ha
averages 150 to 200.

``That's average among my friends,'' she said defensively, glancing at her mother sitting nearby.
``Some send as many as 500 a day.''

In 2009, Ms.~Ha won the South Korean national title against 2.8 million competitors by thumbing 7.25
characters a second. (The best score among participants in their 40s was 2.2 characters a second.)
Mr.~Bae, who was the 2008 national champion and has typed as many as 8 characters a second, did not
compete last year.

``I text while walking, eating, watching TV,'' Ms.~Ha said.

During the interview, even though she was not holding her cellphone, she sat with her thumbs facing
each other like a crab's claws, as if ready to alight on an imaginary keypad.

``At school, we look and listen to the teacher while texting on our cellphones under our desks or in
our pockets,'' she said. ``No typos.''

That behavior has gotten Ms.~Ha, Mr.~Bae and numerous other teenagers around the world into trouble,
with angry teachers confiscating their cellphones. But Ms.~Ha's international victory has more than
compensated for all that.

Ms.~Ha is an almost accidental champion. Last October, she and her friends were walking through the
Coex Mall, a youth hub in southern Seoul, when they saw an LG kiosk about the national competition
that was under way. She entered, hoping to get free movie tickets and other gift coupons. She ended
up with the top prize: \$17,000.

Mr.~Bae was more calculating. ``When I saw the 2008 competition announced on the Internet, I said,
This is it,'' he said. ``It was about time that someone organized something like this. I wanted to
make money and buy a car.''

For the international championship, he trained by copying billboards and anything else that came
into view. He also transcribed passages from the Bible during his flight to New York. Ms.~Ha trained
by typing titles on her bookshelf for five minutes every night before going to bed.

``The more you text, the faster you get,'' she said. ``In the competition, it's not just about the
speed, though. You have to be calm and not make mistakes.''

Ms.~Ha began using a cellphone in the fourth grade, and Mr.~Bae began in the fifth grade. They are
fervent believers in texting. They call it a far more efficient means of communication than a
conversation, a telephone call or an e-mail message.

``You would rather text than dial and wait for the other person to answer,'' Mr.~Bae said. ``It's
especially good when you say 'hi' to someone you haven't seen for a while or don't know well. You
avoid the awkwardness you might feel on the phone. Texting is the modern letter, but I admit it's
not the same as talking face to face.''

Ms.~Ha added: ``When you talk, you often blabber. If you text, you think more coherently because you
have to make yourself understood in short but logical sentences.''

Her mother, Kim Young-sook, 46, said she hoped that Ms.~Ha's world title would help her win
admission to a college as an engineering major. But she still hates it when Ms.~Ha texts while
eating or studying.

``You should show some respect for the food you are eating and the person you are dining with,'' she
said, scolding her daughter. ``Kids with cellphones don't have manners and look so distracted.

``The cellphone is a great gift of civilization but also one of its pollutants,'' Ms.~Kim added.
Shortly afterward, though, she was texting herself.

\section{Obama to Party: Don't `Run for the Hills'}

\lettrine{P}{resident}\mycalendar{Jan.'10}{28} Obama vowed Wednesday night not to give up on his
ambitious legislative agenda, using his first State of the Union address to chastise Republicans for
working in lock-step against him and to warn Democrats to stiffen their political spines.

\myphoto{Obama.to.Party.Don.t.Run.for.the.Hills..jpg}{President Obama delivered his State of the
  Union speech on Wednesday night.}

Mr.~Obama appealed for an end to the ``tired old battles'' that have divided the country and stalled
his efforts on Capitol Hill. He promised to focus intently on the issue of most immediate concern to
the nation, jobs. And with his top priority, a health care overhaul, delayed in the wake of the
recent Republican Senate victory in Massachusetts, he offered a pointed message to both parties.

``To Democrats, I would remind you that we still have the largest majority in decades, and the
people expect us to solve some problems, not run for the hills,'' Mr.~Obama said in his nationally
televised speech. ``And if the Republican leadership is going to insist that 60 votes in the Senate
are required to do any business at all in this town -- a supermajority -- then the responsibility to
govern is now yours as well. Just saying no to everything may be good short-term politics, but it's
not leadership.''

The speech, Mr.~Obama's third to a joint session of Congress, comes at a particularly rocky point in
his presidency, with many Americans -- including some fellow Democrats -- complaining that the
president has lost sight of the priorities of ordinary people. And Mr.~Obama acknowledged their
doubts, conceding that some of his political setbacks ``were deserved,'' a striking admission for
any president.

His tone was colloquial, even relaxed; at one point he joked that the bank bailout was ``about as
popular as root canal.'' But at the same time Mr.~Obama struck a defensive note, reminding the
nation yet again that he inherited a mountain of problems and insisting that, one year after he took
office, ``the worst of the storm has passed.''

At a time when many Americans are concerned, even angry, about the economy and about the performance
of government more generally, Mr.~Obama sought to restore public confidence in his administration
and to persuade Americans that he is directing his attention more fully to the economy. While he did
not offer any sweeping new agenda or far-reaching legislative program, he put forth a handful of new
initiatives, including plans to provide small businesses with tax breaks and better access to bank
loans.

After refusing to set a timetable for the repeal ``Don't Ask, Don't Tell,'' the military's policy
barring openly gay men and lesbians from serving, he vowed to work with Congress this year to repeal
it. He called for the reauthorization of No Child Left Behind, his predecessor's signature education
law. In a nod to the growing political and economic pressure to begin reining in the budget deficit,
he proposed a freeze on a portion of the domestic budget.

Mr.~Obama campaigned on a promise to change the culture of Washington and to make government
transparent. But on Wednesday night, he suggested that he believed he had not done enough, and spoke
of a ``credibility gap'' that must be closed by curbing the outsized influence of lobbyists. ``We
have to recognize that we face more than a deficit of dollars right now,'' he said.

Reprising a line he used in last year's address to Congress, he said, ``We face a deficit of trust
-- deep and corrosive doubts about how Washington works that have been growing for years.''

He called for new rules requiring lobbyists to disclose each contact they make on behalf of a client
with Congress or with his administration. And, in a rare flash of open confrontation between the
White House and the Supreme Court, Mr.~Obama declared that a recent court ruling would ``open the
floodgates for special interests,'' and perhaps foreign companies, to exert more influence in
political campaigns. Justice Samuel A.~Alito Jr., breaking with decorum at such events, shook his
head and appeared to mouth the words, ``No, it's not true.''

Republicans said they welcomed the president's partial freeze on domestic spending. But they warned
against what they regard as the president's big government agenda. In delivering his party's
response, Gov. Robert F.~McDonnell of Virginia, a newly elected Republican, declared, ``The
circumstances of our time demand that we reconsider and restore the proper, limited role of
government at every level.''

But rather than retreat from his ambitious agenda, Mr.~Obama sought Wednesday to repackage it, by
explaining how his top priorities -- the health measure, tough new regulations on banks, energy
legislation -- fit into his broader initiative to put the economy on sounder footing for the long
run.

On health care, Mr.~Obama did not chart a specific path forward for Congress. Rather, he appealed to
lawmakers to ``take another look at the plan we've proposed'' once temperatures cool after the
Republican win in the Massachusetts Senate race. He added, ``Do not walk away from reform. Not now.
Not when we are so close. Let's find a way to come together and finish the job for the American
people.''

Still, after a year of working to get health care passed, Mr.~Obama said his No.~1 issue is now the
economy and jobs. ``Jobs must be our No.~1 focus in 2010,'' Mr.~Obama said, adding ``People are out
of work. They are hurting. They need our help.''

To that end, the president renewed his call for Congress to pass a jobs bill that would spur
investment in green jobs and clean energy, though he did not offer specifics of what it would cost.
He proposed investment tax cuts that would put more cash in the pockets of small business owners and
a new program that would take \$30 billion from the fund used to bail out troubled banks and
automakers, and redirect it toward an initiative to encourage community banks to lend to small
businesses.

He set a goal of doubling exports over the next five years -- an increase that he said would support
two million jobs. And, as he pledged to do earlier in the week, Mr.~Obama also outlined a series of
proposals intended to help the middle class, including new tax credits for child care and a cap on
student loan payments for recent graduates.

And Mr.~Obama offered a very public show of confidence in one of the architects of his economic
plan: Treasury Secretary Timothy F.~Geithner, whose close identification with Wall Street has made
him a focus of some of the populist anger directed at the White House. When Mr.~Obama strode into
the chamber of the House of Representatives to deliver the address, he stopped to face Mr.~Geithner,
who had just spent the day getting grilled on Capitol Hill and put both hands encouragingly on the
secretary's shoulders.

Strikingly, for a president who is prosecuting two wars and trying to protect the country against
the threat of a terrorist attack, Mr.~Obama spent only nine minutes in an address that lasted more
than an hour on foreign policy. He renewed one of the most popular promises of his campaign for
election, to bring the troops home for Iraq, saying ``Make no mistake -- this war is ending, and all
of our troops are coming home.''

But he devoted only one paragraph to a far less popular decision, escalating the troop levels in
Afghanistan. ``There will be difficult days ahead,'' Mr.~Obama said. ``But I am confident we will
succeed.''

As have presidents before him, Mr.~Obama grappled with how to describe the state of the union. In
the end, he settled on the formulation that many of his predecessors have used, with a twist:
``Despite our hardships, our union is strong.''

\section{The Afghan Leader's Hat, Always More Than Just Headgear, Is Losing Its Cachet}

\lettrine{I}{t}\mycalendar{Jan.'10}{28} once attracted the admiration of trendsetters in the West,
the jests of comedians at home and abroad, and the somewhat impotent ire of animal rights advocates.

Hamid Karzai's hat, while still firmly on the Afghan president's head whenever he appears in public,
is no longer quite the symbol it once was.

Known as a karakul hat, and made of the pelt of fetal or newborn lambs of the karakul breed of
sheep, traditionally it was something worn by Tajiks and Uzbeks from northern Afghanistan. When
Mr.~Karzai, a Pashtun from the turban-wearing south, took office in 2002, the karakul hat was part
of his attempt to devise a wardrobe that was Afghan rather than ethnic or regional.

\myphoto{The.Afghan.Leader.s.Hat.Always.More.Than.Just.Headgear.Is.Losing.Its.Cachet.jpg}{President
  Hamid Karzai of Afghanistan at a news conference in Kabul last December.}

It was a move widely praised at the time, in Afghanistan and abroad. The American designer Tom Ford
called the Afghan president ``the chicest man on the planet.'' Afghans looking for national symbols
after decades of ethnic strife inspired a brisk trade in the hats, made of lambskins from
Mazar-i-Sharif in the north and fashioned by Kabul's hatters, whose shops lined both sides of
Shah-e-do Shamshera Wali Road.

Now, a tainted presidential election later, and with efforts to make a truly multiethnic government
foundering, the sheen is off the shimmery fur headwear.

Young men no longer wear it; Mr.~Karzai's opponent in the aborted election runoff, Abdullah
Abdullah, a northerner, preferred a hatless suit-and-tie ensemble. All but 12 of the hatters shops
have closed on Shamshera Road, also famous for its shrine covered in pigeons. Those remaining say
they are lucky to sell a hat a day.

``I went back to my village in Logar wearing my karakul hat,'' said Ahmed, an Afghan in his 50s, who
was shopping for a new hat, ``and people laughed: 'There goes the old man who thinks he's
president.' '' It was not clear which offended him more, ``old'' or ``president.''

``Hamid may be the only guy in Afghanistan wearing that particular kind of hat,'' said a post on a
satirical Web site, Ridiculopathy.com, ``but all the same the pointy wool chapeau has come to
symbolize the country to the rest of the world.''

Just as Mr.~Karzai's hat is more than just a hat, the reaction against it is more than just a
fashion whim. ``It would have been better if he just wore a turban. It would have been more
honest,'' said Rahnaward Zariab, a novelist and cultural commentator on Tolo TV in Kabul. ``Instead
he deceived the nation. The costume of Karzai doesn't mean anything; it's not a symbol anymore. Now
we are seeing his actions, and it's clear now that he is a Pashtun.''

Mr.~Zariab complained that there were relatively few non-Pashtuns in Mr.~Karzai's new cabinet, which
is yet to be completely approved by Parliament.

Efforts to solicit a comment from the president on his headgear met with no success, and slight
annoyance. ``Everything else is finished with,'' said his spokesman, Waheed Omer, ``Now you're going
to write about the hat?''

Mr.~Karzai himself once, at a military ceremony in Kabul, explained his affection for the karakul
hats. ``I wear them because they are very, very Afghan,'' he said, according to an Associated Press
account. ``And if it looks good, all the better.''

Among the hatters, at least, the president still gets rave reviews for his good taste. He is also
one of their best customers.

Mr.~Karzai's affection for the karakul hat is so strong that, if the hatters' accounts are to be
believed, he has purchased dozens of them since taking office. Sayed Habib Sadat, owner of one of
the remaining hat shops near the shrine, says he has sold Mr.~Karzai 15 karakul hats in various
shades, mostly the dark gray he prefers, but also blacks, and mottled brown-and-whites.

``The president is bringing the old traditions together and showing people: 'I'm an Afghan. I'm
using my own tradition,' '' Mr.~Sadat said. ``It's been good for him, and for us.''

Karakul hats are not cheap; good ones sell for hundreds of dollars, and some can reportedly cost up
to \$3,000.

The more expensive ones are made from the skins of lambs taken from the pregnant ewe just before
birth, by cutting open her abdomen, sometimes while she is still alive. Less costly are those made
from lambs killed immediately after delivery; because karakul sheep are extremely protective of
their young, that often means slaughtering both together, or forcibly separating them.

The slaughterer first palpitates the belly of the ewe, feeling the hooves of the fetus; if they are
hard, that means a particularly good lambskin if it is harvested pre-partum.

A single lambskin makes one hat. First the skin is sewn into a cylinder. Then it is stretched over a
head-shaped mold, or qalib, carved from a chunk of wood from a mulberry tree. Finally, the hats are
pounded into their rectangular shape with a rounded river stone.

The more expensive are particularly soft and flat, with the texture of moir\'e silk and a velvety
nap, arranged in fine, patterned ridges. Sometimes those ridges seem to spell out the name of Allah
in Arabic, which makes them particularly valuable. On rare occasions they may even appear to spell
out, in Arabic, the Kalimah or Muslim declaration of faith -- ``There is no God but God and Muhammad
is his prophet'' -- making them even more dear.

That does not, however, endear them to the religious authorities. Because the sheep is not killed in
the way prescribed by Islam, the hat and even the meat of the sheep from which it is made are haram,
or forbidden.

When Mr.~Karzai orders a hat, the hatters are summoned to the Presidential Palace to measure his
head personally. ``When I first made him a hat, he was size 22 ½ inches,'' Mr.~Sadat said. ``Now
he's 23 ½ inches.'' Other hatters reported similar measurements, give or take half an inch.

Mr.~Sadat said Mr.~Karzai was not unique in this regard. He recalled Noor Mohammed Tarakai,
Afghanistan's first communist president in 1978, who also went from a 22 inch hat size to 23 inches,
and that was during a much shorter term in office; he was assassinated after a year, reportedly by
smothering.

``Presidents get swollen heads,'' Mr.~Sadat said.

\section{Economic Measures by North Korea Prompt New Hardships and Unrest}

\lettrine{N}{orth}\mycalendar{Feb.'10}{04} Korea's sweeping attempts to revive socialist central
economic planning and crack down on private markets appear to have set off runaway inflation and
perhaps caused the country's reclusive leaders to reconsider some of their own measures, according
to South Korean news reports and intelligence officials.

\myphoto{Economic.Measures.by.North.Korea.Prompt.New.Hardships.and.Unrest.jpg}{Old North Korean
  currency bills, center, were on sale with the currencies of other countries at a market in Beijing
  last month.}

As one indicator of a possible shift in the leadership's approach, North Korea's leader, Kim
Jong-il, is reported to have dismissed the official who oversaw a sweeping monetary conversion late
last year that led to severe price increases and social unrest, which is unusual in North Korea, a
tightly policed nation.

The official, Pak Nam-gi, chief of the finance and planning department of the ruling Workers' Party
of North Korea, was blamed for skyrocketing food prices that had severely curtailed the country's
already meager food supplies, according to the Chosun Ilbo, a mass-circulation daily newspaper in
Seoul.

The Chosun Ilbo article, which was published Wednesday, cited unidentified diplomatic sources in
China. The South Korean government could not confirm the report. Mr.~Kim moves top party officials
so secretively that some of those reported to have been banished have often later resurfaced with
senior posts.

But analyses by South Korean experts who monitor the situation across the border suggest that a
series of steps by Mr.~Kim's government to close private markets, radically revalue the currency and
restrict the importing of consumers goods from China have created a new crisis for the country's
perennially ailing economy.

Prices of basic goods have sharply risen since the government replaced the old bank notes in late
November, exchanging them for new bills at a rate of 100 to 1, according to officials here and
reports on Seoul-based Web sites that glean information from sources inside the North. North Korea
placed a cap on how much old money a citizen could exchange for the new currency, effectively
confiscating private stores of wealth. Reports said it also outlawed foreign currency holdings and
punished private traders who smuggled goods from China.

But in an effort to reward people who depend on state salaries, like factory workers and members of
the military, the government paid the same nominal wages in new currency, theoretically increasing
their purchasing power a hundredfold. The underlying idea was to suppress the country's black
markets and reassert the economic model of central planning.

``The problem with the currency reform was that it hurt imports through unofficial channels while
state-run stores failed to stock their shelves with goods,'' said Dong Yong-sueng, an expert on the
North Korean economy at the Samsung Economic Research Institute in Seoul. ``The factory workers,
with their boosted purchasing power, turned to the unofficial markets, driving up the prices.

``Whether this is a temporary problem or the regime will eventually take control of the situation --
we won't find out until one or two months from now.''

The South Korean unification minister, Hyun In-taek, also said it was too early to measure the
impact of the currency changes in the North.

Illegal markets had been mushrooming in North Korea since the government's ration system failed to
supply enough food, leading to a devastating famine in the mid-1990s. A new class of traders
emerged, operating outside the official system and establishing corrupt ties with party officials,
according to analysts in Seoul.

The emergence of a merchant class and repeated visits by Mr.~Kim to different cities in China were
taken as signs that the North was cautiously embracing elements of a market economy. But the steps
taken in the past few years suggest a sharp shift in the other direction.

Mr.~Pak, the official said to have been dismissed, took the top financial job in 2005 as North Korea
began introducing a series of measures aimed at suppressing the black markets. He led the currency
reform, which North Korea hoped would arrest inflation by eliminating profiteering in private
markets. That was viewed as an important task to be carried out before the youngest of Mr.~Kim's
three sons, Kim Jong-un, was to be officially anointed as the leader's successor, around 2012, the
Chosun Ilbo said.

But the currency revaluation instead fueled much higher inflation and uncertainty, bringing market
activities nearly to a halt and driving people to panic, South Korean reports and analysts said.
Mr.~Pak disappeared from reports in North Korea's state-run media in early January.

A group of veterans from the 1950-53 Korean War staged a protest last month, sitting in front of a
party office in Danchon, a town on North Korea's eastern coast, according to Good Friends, a relief
group based in Seoul. The unusual show of public discontent in North Korea took place after some
people in Danchon died of hunger, the group said. The government quickly released emergency food in
late January to quell the unrest, it added.

``As the food situation worsened in the North,'' Good Friends said, ``there is talk there that if
the government did not come out with drastic measures to improve the situation by the end of
February, it could trigger an uprising.'' Good Friends runs one of several Web sites that provide
news on North Korea, using sources inside the North. Its reports often describe small sectors of
North Korean society.

In some parts of North Korea, food prices have more than tripled in the past three weeks, according
to Daily NK, another Web site that specializes in North Korean news. The value of dollars on the
black market, it said, was soaring as well.

The inflationary surge appears to have two related causes. One is that private sales of goods have
become far riskier, pushing up the price of any goods available outside the state distribution
system.

The other is that the currency move at least temporarily rewarded the elite and the loyal class of
state workers, who received sharp salary increases they could spend at state stores. But the North's
centrally planned economy has long lacked the ability to produce enough food and consumer goods at
state-set prices, meaning demand outstripped supply even more than it usually does.

\section{Reactionaries? Make That `Collectors'}

\lettrine{C}{contestant}\mycalendar{Feb.'10}{04} No.~3, a portly man in suspenders named Cui
Xiaosong, clutched a golden mallet and gulped like an executioner having second thoughts. As a guest
on China's wildly popular antiques reality show ``Collection World,'' Mr.~Cui knew he might have to
get violent before the next commercial break. The victim? A delicately painted vase he had brought
to the show, which he believed to be from the Qing dynasty and worth about \$30,000.

\myphoto{Reactionaries.Make.That.Collectors.jpg}{A worket at the ACF China furniture factory with a refurbished trunk.}

``If it's a fake, will you smash it?'' asked the program's white-gloved host, Wang Gang, as Mr.~Cui
faced the studio audience and three guest judges.

Mr.~Cui nodded. The audience quieted down and Mr.~Wang used the final minute to impart a bit of
wisdom about collecting antiques in modern-day China: ``Just as China opened up, so too is
collecting about opening the mind to understand the outside world.''


It was hard to tell whether Mr.~Cui was listening, but he certainly heard the host announce the
judges' verdict: ``It's a modern reproduction!''

Mr.~Cui winced as he swung the mallet, shattering the vase -- and with it his dreams of the wealth
it might have brought at auction. Cue the instant replay.

Some four decades after the Cultural Revolution, when many of the country's centuries-old treasures
were defaced or destroyed as a result of Mao's command to eradicate ``the four olds'' -- old ideas,
old culture, old customs and old habits -- China has reversed its attitude toward antiques. Ming
dynasty porcelain vases, 19th-century hardwood furniture and even early 20th-century calligraphy ink
pots have become popular status symbols for an emerging middle class eager to display its new wealth
and cultural knowledge. The antiques market has become so hot, in fact, that it has given rise to a
new category of must-see TV here.

In recent years, ``Collection World'' and a dozen other similar shows -- with names like ``Treasure
Appraisal'' and ``Art Collector'' -- have been luring both serious collectors and armchair
enthusiasts, offering information on collecting trends and appraisal techniques, and encouraging a
new wave of treasure hunting.

While some in the antiques world laud these programs for turning antiquing into a national pastime,
others are skeptical of their educational value. As Yan Zhentang, the president of the Chinese
Collectors' Association, noted, ``These shows certainly help get ordinary people interested in
antiques, but the bottom line is they are just entertainment, and they make mistakes.''

\myphoto{Reactionaries.Make.That.Collectors.1.jpg}{Shopping the Panjiayuan antiques market.}

Daniel Newham, a British expatriate who has become a popular television personality in China, said
he was dismayed by the lack of professionalism when he served as a celebrity judge on an episode of
``Collection World.''

``The other judges were pretty awful,'' Mr.~Newham said, adding that one of them admitted to him
that he had only recently started working in the field of antiques and did not have the skills to
properly appraise the featured items. (The show's executives declined to comment and refused to
allow Mr.~Wang, the host, to be interviewed.)

Nevertheless, the shows have attracted a devoted following. Zhou Yajun, a long-distance truck driver
and collector from Hebei Province, near Beijing, said he watched ``Collection World'' and other
antiques shows every week, testing his appraisal skills against those of the judges in the hope that
he could learn to outwit the counterfeiters who prey on the country's amateur antiquarians.

Mr.~Zhou, 38, said he began collecting antiques four years ago, and his hobby quickly became
all-consuming. ``For a week after I bought my first antique, I would hug it to sleep, I was so
excited,'' he said, showing off photos of his favorite purchases on his cellphone during a morning
of poking around Panjiayuan, Beijing's vast antiques market.

Mr.~Zhou said he had spent the equivalent of \$12,000 so far feeding his addiction, a hefty sum for
a man who earns less than \$18,000 a year. But spending so much time alone on the road takes an
emotional toll, and collecting has become a way to fill the void.

``If I don't see my antiques for a few days, I miss them,'' he said.

``The problem is, everyone wants to collect now, so there's not much of the real stuff left,'' he
added, eyeing some rusty coins advertised as 100 years old before shaking his head and moving on to
the next vendor.

Distinguishing real Chinese relics from their latter-day replicas can be a daunting task, especially
since forgers have access to the same televised information that collectors do. ``I used to go to
the countryside to buy antiques,'' Mr.~Zhou said. ``But lately I've found the peasants are buying
fakes and making up a story to pass the pieces off as authentic.''

Perhaps wisely, Mr.~Zhou has come up with his own way of evaluating authenticity: ``After I buy
something, I put it in my home for two days,'' he said. ``If I start to like it, it's real. If not,
it's counterfeit.''

THE Chinese government has become increasingly assertive about claiming ownership of its national
heirlooms. It condemned Christie's last year for auctioning bronze sculptures looted from the
capital's Old Summer Palace in 1860 and, more recently, it sent out government officials and art
historians to inspect the collections of global art institutions like the Metropolitan Museum of Art
and the Smithsonian for cultural artifacts that might have been illegally obtained from China. And
so, many private collectors have come to regard their passion not just as a smart investment, but as
a patriotic duty.

``Chinese people are becoming richer and need to be responsible for our dignity and history,'' said
He Shuzhong, the deputy director of the State Administration of Cultural Heritage's legal and policy
department and the founder of the Beijing Cultural Heritage Protection Center, a nonprofit
organization. ``How can China rise peacefully if we cannot protect our culture?''

But many in the industry acknowledge that the profits driving the antiques trade are a more powerful
incentive than nationalism.

As Yan Xubao, 31, a dealer at the ACF China furniture company in the Gaobeidian market on the
outskirts of Beijing, observed, ``Without a free capitalistic spirit, these antiques would still be
buried in the countryside somewhere.''

Mr.~Yan is a regular at many of the city's antiques wholesale markets, where peasants bring old
broken furniture, farming tools and stone carvings collected from the outer provinces. Such items
are bought by urban restorers, like those at ACF, who resell the repaired pieces, often at a huge
markup.

While the global economic crisis has affected ACF's wholesale business, which often exports to
retailers abroad, its retail sales have remained relatively robust because of the strength of the
Chinese economy and the antiques industry's growing grassroots base in China, said Roger
Schwendeman, founder and one of the company's managing partners.

Mr.~Schwendeman, an American who has worked in China's antiques trade for eight years, said Chinese
buyers are still paying top dollar for jade and furniture from the Ming and Qing dynasties made from
rare hardwoods like yellow rosewood and ebony, which most foreigners ignore.

``Western buyers ask about history, while Chinese are interested in the value of the material,'' he
said, over the noise of hammering and sawing, as a trio of workers restored an ornately carved
rosewood cabinet at his factory outside Beijing.

Many of those same foreigners who bought up troves of China's antiques in the 1980s and '90s are now
seeking out the increasingly wealthy mainland Chinese buyers, Mr.~Schwendeman added. ``They know the
money and passion are in China.''

\section{At Least 5 Dead, 12 Hurt in Connecticut Power Plant Blast}

\lettrine{A}{}\mycalendar{Feb.'10}{08} power plant under construction in central Connecticut
exploded with earthquake force that shook homes across much of the state on Sunday as workers purged
natural gas lines in preparation for the plant to open this year. At least five people were killed
and more than two dozen were injured as a section of the plant collapsed and burned.

\myphoto{At.Least.5.Dead.12.Hurt.in.Connecticut.Power.Plant.Blast.jpg}{Smoke rose from the the Kleen
  Energy plant on River Road in Middletown, Conn., on Sunday.}

Witnesses said the explosion at the Kleen Energy Systems plant in Middletown, Conn., 15 miles south
of Hartford, occurred at 11:17 a.m. in a thundering convulsion of flames and smoke seen for miles
around and felt as far away as cities and towns on the shore of Long Island Sound, 30 miles away.

As towering plumes of dark smoke poured into a dazzling blue sky, scores of ambulances, fire
engines, police cars and helicopters streamed to the scene on the west bank of the Connecticut River
on the southern outskirts of the city, the home of Wesleyan University.

Fire and rescue teams from Middletown, Durham, Portland, Cromwell and other towns converged as crews
fought the blaze into the afternoon. Aerial pictures showed a smoking, sprawling riverside site with
buildings housing generators, fuel tanks and other power equipment, topped by two smokestacks. The
explosion apparently occurred in the largest building, called the Power Block, which was destroyed.

River Road leading to the site was a tangle of firefighting and rescue equipment. Flames were seen
shooting from a pipeline after the blast, but the line was capped shortly after noon, officials
said. Even so, scattered fires blazed and smoke billowed over the scene for hours, and
search-and-rescue teams looked through the rubble for victims into the night. The search, with dogs,
sound-detecting and thermal imaging equipment, could go on for days, officials said.

Mayor Sebastian N.~Giuliano of Middletown, at a late-afternoon news briefing, said that five people
were known to be dead, but he did not release names pending notification of their families. The son
of one victim, Raymond Dobratz, 57, of Old Saybrook, a pipe fitter who had been working at the plant
for a year, said his father had been killed.

Mayor Giuliano said the number of dead and injured remained unclear because the number of workers at
the plant fluctuated, in part because so many subcontractors were involved. He said the workers had
been purging gas lines all weekend.

He said that as many as 200 construction workers had been at the plant daily, though fire officials
said only about 50 were on the job at the time of the explosion. The mayor said various contractors
were being asked to identify workers at the site on Sunday. ``We need to know who was there today,''
the mayor said.

And while the cause of the blast remained undetermined, the mayor said that a natural gas explosion
was ``the assumed cause.'' He added, ``Terrorism has been ruled out.''

Mr.~Giuliano said that many of the construction workers had been evacuated from the site before the
line purging, as standard procedure, and he said there had been no previous accidents at the plant
-- ``not so much as a hangnail,'' he said.

The explosion was under investigation by the state and local police and by several federal agencies,
including the Occupational Safety and Health Administration, the Justice Department's Bureau of
Alcohol, Tobacco, Firearms and Explosives and the United States Chemical Safety Board, which
examines chemical industrial accidents. Daniel M.~Horowitz, a spokesman for the chemical safety
board, said that a team of investigators would be at the site on Monday.

Fire officials said the explosion occurred as workers for the construction company, O \& G
Industries, were purging the pipelines of natural gas, the main source of fuel for the plant, in a
procedure known as a blow-down.

Mr.~Horowitz said he could not confirm that report, but noted that a gas purge last June at a
food-processing plant in Garner, N.C., killed four people and injured 67 others. He said his agency
had issued safety recommendations in that case just last week.

Gordon Holk, the Middletown plant's general manager, said workers from O \& G; Ducci Electrical
Contractors, of Torrington; and Keystone Mechanical Electrical Contractors had been at the site on
Sunday.

Tests were under way in preparation for a spring or summer opening of the 620-megawatt plant, which
has been under construction since September 2007 in an old feldspar quarry and was 95 percent
finished. The plant, one of the largest power facilities to be built in New England in recent years,
was to supply electricity to Connecticut Light and Power.

As many as 1,000 workers had been employed in building the plant, but the number had declined to 400
or 500 recently, according to Philip Armetta, the project's developer, whose 45 percent interest was
placed in a trust when he pleaded guilty in 2007 to a charge of failing to disclose knowledge of
price-fixing in the state's trash disposal industry.

Representative Matthew Lesser, a freshman Democrat who lives a mile from the plant and represents an
adjacent district in the State House, said the \$1 billion plant -- its costs raised by delays due
to what he called ``regulatory hurdles'' -- had been expected to cut the costs of power in a state
that has some of the highest rates in the nation.

Al Santostefano, the deputy fire marshal in Middletown, said the authorities believed that about 50
employees were in the building where the explosion occurred, and he spoke of the difficulties of
searching for victims in the rubble. ``It's a slow dig,'' he said. ``There's a possibility someone
could still be alive under the debris.''

While the number of casualties was uncertain, hospitals in Middletown, Hartford and New Haven
reported receiving more than two dozen injured people.

At Middlesex Hospital in Middletown, Peg Arico, a spokeswoman, said 26 people were brought from the
site for treatment. Two victims suffering major injuries that could not be treated at Middlesex were
transferred to trauma centers, one to Hartford Hospital and one to Yale-New Haven, Ms.~Arico said.
Most of the others were treated and released.

``We did not see a lot of burn victims,'' she said. ``Of the ones that we are still treating, many
have broken bones and others have abdominal pain. The injuries can be described as impact injuries
from the explosion.''

Hartford Hospital said two injured people were brought directly from the explosion, in addition to
one transferred from Middlesex .

The explosion shook the walls and windows of homes 20 to 30 miles from the scene and touched off an
avalanche of telephone calls as residents across a swath of the state contacted relatives and
friends to see if anyone had been hurt. It also set off a flood of Internet commentaries.

The project contractor, O \& G Industries, was described on its Web site as a closely held company
based in Torrington, Conn. Algonquin Gas Transmission was listed as the gas supplier, and the
plant's turbines were manufactured by Siemens Power Generation. Kleen Energy Systems is owned
principally by Energy Investor Funds, which recently acquired an 80 percent share.

Mr.~Lesser, a member of the Energy and Technology Committee of the State House, said he was having a
cup of coffee when his building shook.

``There was a loud rumble and my windows in my apartment rattled for 5 to 10 seconds,'' he recalled.
``I had no idea what it was. It was peculiar, but I didn't think anything of it.'' Then, he said, he
began getting phone calls and text messages about the explosion.

``The first couple asked, 'Did you feel that?' and one person thought it was an earthquake,'' he
said. ``And then the subsequent ones reported about the explosion itself.''

The tests being conducted at the plant on Sunday were in preparation for a spring opening, he said.
``The hope was that by increasing generation, we could bring electric rates under control, which in
the state of Connecticut, it is my understanding, are the highest in the continental United
States,'' he added.

\section{Obama Meets Dalai Lama, and China Is Quick to Protest}

\lettrine{P}{resident}\mycalendar{Feb.'10}{23} Obama met with the Dalai Lama on Thursday, welcoming
the exiled Tibetan spiritual leader to the White House for a low-profile meeting that nonetheless
raised the hackles of China.

\myphoto{Obama.Meets.Dalai.Lama.and.China.Is.Quick.to.Protest.jpg}{The Dalai Lama spoke to the news
  media outside the White House after meeting with President Obama on Thursday.}

The two men spoke about democracy, human rights and the need to preserve Tibet's religious identity
and culture -- all issues that, predictably, irritated Beijing.

In a written statement, the White House said Mr.~Obama had expressed support for the preservation of
Tibet's ``unique religious, cultural and linguistic identity and the protection of human rights for
Tibetans in the People's Republic of China.''

The meeting, which the White House put off last year so as not to interfere with Mr.~Obama's trip to
Beijing in November, is the latest evidence of the more muscular posture that the Obama
administration has been adopting with China, after a year in which China has taken a firm stance in
its relations with the United States on a number of issues.

Last month, the administration announced a \$6 billion arms sales package to Taiwan, infuriating
Chinese officials.

The meeting on Thursday between Mr.~Obama and Tenzin Gyatso, the 14th Dalai Lama, did not please
Chinese officials, either.

China, which regards the Dalai Lama as an advocate of Tibetan independence, said that it was
``strongly dissatisfied'' and that it expected the United States to try to make amends.

On Friday, China accused Mr.~Obama of ``seriously damaging'' ties between the two countries, the
spokesman for the Chinese Foreign Ministry, Ma Zhaoxu, wrote in a statement posted on the ministry's
Web site.

China also summoned the American ambassador, Jon Huntsman, to lodge a formal complaint about the
meeting at the White House, the statement online said.

While the White House took pains to avoid the appearance that this was a meeting between heads of
state -- it took place in the White House Map Room, not in the Oval Office -- the Dalai Lama also
met with Secretary of State Hillary Rodham Clinton at the State Department.

``The Dalai Lama is a Nobel Peace Prize laureate, internationally revered religious and cultural
leader,'' said Mark Toner, a State Department spokesman, explaining why the country's senior
diplomat would meet with him. Mr.~Toner said other secretaries of state had also met with him.

The White House went out of its way to keep the meeting low-key.

White House officials did not allow reporters or photographers to see Mr.~Obama and the Dalai Lama
together, and they released a single official photograph after the meeting.

The Dalai Lama did meet with reporters at the White House afterward, and said he was happy with the
visit. He also tossed some snow at reporters.

In deference to China, American presidents usually do not meet publicly with the Dalai Lama.
President George W.~Bush broke with that tradition in 2007 when he attended a ceremony in the
Capitol Rotunda in which Congress gave the Dalai Lama its highest civilian honor, the Congressional
Gold Medal.

The Dalai Lama has lived in exile in India since the Chinese Army crushed an uprising in his
homeland in 1959, and he is revered as the spiritual leader of Tibetan Buddhists. He has been
pressing to return to China to advocate for greater cultural and religious freedoms for his
followers. He has said he wants only genuine autonomy for Tibet, not independence.

\section{N.H.L.'s Stance on Olympics Faces Challenge}

\lettrine{G}{ary}\mycalendar{Feb.'10}{24} Bettman could be finding himself boxed into a corner, his
position that N.H.L. participation in future Olympics may not be worth the sacrifice rendered
untenable by the big television audience last Sunday for the game between Canada and the United
States.

Bettman, the N.H.L. commissioner, has suggested that shutting down for 15 days may be more costly to
the league than the benefits it derives from having its stars showcased internationally. The league
gets no revenue from Olympic telecasts, but the high viewership numbers suggest that the N.H.L. is
getting more exposure from the Olympics than Bettman or anyone else expected.

\myphoto{N.H.L.s.Stance.on.Olympics.Faces.Challenge.jpg}{Ryan Miller, goaltender for the Buffalo
Sabres, during Sunday's game against Canada.}

Even if his reluctance to embrace participation in future Winter Games is mainly a bid for leverage
in negotiations for a piece of the Olympic pie or in talks with the N.H.L. Players' for a new
collective bargaining agreement, it is becoming increasingly difficult to envision saying goodbye to
a platform that provides undeniable benefits to the N.H.L.

According to figures, the game had 10.8 million viewers in Canada, making it the most-watched sports
program in Canadian television history, and 8.2 million in the United States on the cable news
channel MSNBC, matching the channel's election-night viewership.

It was not even an elimination game, yet viewership in both countries surpassed that of Game 7 of
last spring's Stanley Cup finals, which was carried in the United States on the regular NBC network.

The numbers for the game Sunday seem to contradict the notion that the Olympics are a too-costly
sacrifice for his league. Rather, the presence at the Olympics of N.H.L. players seems to be putting
the spotlight on the league in a way the N.H.L. itself has not quite been able to do.

``It's great for the league and hockey as a whole, a chance to showcase the best players in the
world on a world stage,'' Nicklas Lidstrom, who plays for Sweden and the Detroit Red Wings, said
Monday night.

Joe Sakic, formerly of Canada and the Colorado Avalanche, said: ``The players love doing it, the
fans love it, it grows the sport. Everyone at the end of the day realizes it's the best thing for
hockey.''

Bettman made some persuasive points last Thursday when he listed the reasons the N.H.L., which in
1998 began taking a break for the Olympics, may decide not to send its players to the 2014 Winter
Games in Sochi, Russia.

``It costs us money, it disrupts our season,'' he said about the league's Olympic shutdown. He cited
other issues as well, including the control the International Olympic Committee asserts over rights
and media access, NBC's seemingly scattershot scheduling of Olympic matches on its ancillary
networks and, of course, the risk of injury to its best players.

Among the things Bettman may be seeking from the I.O.C. may be visibility for the N.H.L. trademark
at the Games alongside the International Ice Hockey Federation's, which is currently present, and
the same access to live online coverage on the N.H.L.'s Web site that the international federation's
site now has.

Bettman is not alone among N.H.L. executives in expressing doubts about taking part in the Sochi
Games, which travel and time-zone differences make more difficult than participating in the Games
here.

Even Brian Burke, the general manager of the United States men's team but also the president and
general manager of the Toronto Maple Leafs, has publicly raised doubts about the advisability of
continued N.H.L. participation in the Olympics.

``We suspend our season for more than two weeks, and we get very little benefit from it,'' Burke
said last August at the team's orientation camp.

``We get some bang when the Olympics are in North America, in Vancouver or Salt Lake, but when
they're overseas, we have gotten no bang at all. I wouldn't be surprised if this is the last
Olympics the N.H.L. takes part in.''

It's hard to square those statements with the obvious passion Burke has for the United States men's
team. But as a Leafs executive, he was speaking for a significant faction within the league.

Such views were widespread within league front offices then. Whether they still will be in the
aftermath of what is shaping up as one of the greatest showcases ever for the N.H.L. remains to be
seen.

Certainly the players are for going to Sochi four years from now.

Internal surveys have shown that a majority of the players' union favor participation, which is sure
to be a factor when it negotiates a collective bargaining agreement with the league at the end of
next season.

And Alex Ovechkin, Evgeni Malkin and Ilya Kovalchuk have declared that they will leave their N.H.L.
clubs, if necessary, to take part in the 2014 Games in their homeland.

No other players or international federations have indicated what they will do if the N.H.L. does
not cooperate with the I.O.C. in 2014.

Lidstrom and Sakic spoke after a ceremony here to honor the 22 men who have won the Stanley Cup, an
Olympic gold medal and the international federation's World Championship. All 22 in what the
federation calls the Triple Gold Club were present, including eight playing at these Games:
Lidstrom, Peter Forsberg, Henrik Zetterberg, Niklas Kronwall and Fredrik Modin of Sweden; Jaromir
Jagr of the Czech Republic; and Scott Niedermayer and Chris Pronger of Canada.

Another triple gold winner, the Soviet defenseman Vyacheslav Fetisov, a former Russian minister for
sport and now chairman of the Kontinental Hockey League board of directors, not surprisingly spoke
in favor of the N.H.L. coming to Sochi. He said that if the league skipped 2014, there might be
further consequences.

``When they get out of the Olympic family, I think there's no chance to get back,'' Fetisov said.
``It's not hard to find a couple of weeks to send the best players to sell the game.''

The K.H.L. wants everyone at their Olympics, of course. Like all European leagues, it makes its top
players available for international competitions.

Rene Fasel, president of the international federation, acknowledged Bettman's concerns about media
and branding questions but said there were bigger issues.

``Gary should really understand that the game and the fans want to see the best against the best,''
Fasel said. ``When I see the numbers of TV viewers in Canada and the U.S., they're huge.

``I understand when Gary says, 'We lose money,' but it's not that. Gary should not forget that
people in Asia and Europe are watching this tournament, and that's also promotion for the N.H.L.''

The increased interest in Olympic hockey in North America raises additional issues about its appeal
versus that of the N.H.L. form of the game.

Fighting simply does not happen in Olympic hockey, at this or any other Olympic tournament, and no
one seems to miss it. Nor does anyone seem to miss the N.H.L. All-Star Game, which Bettman cited as
another sacrifice the league makes in Olympic years.

But in the end, the numbers seem to say some definite things that Bettman must keep in mind when it
comes time to decide about Sochi.

Olympic hockey creates more interest than N.H.L. hockey, but Olympic hockey without N.H.L. players
in it would create only a fraction of that interest.

\section{Transplants That Do Their Job, Then Fade Away}

\lettrine{J}{onathan}\mycalendar{Feb.'10}{24} Nu\~{n}ez was 8 months old when a liver transplant
saved his life. Three years later, his body rejected the transplant, attacking it so fiercely that
it wasted away and vanished, leaving barely a trace.

That result, seemingly a disaster, was exactly what his doctors had hoped for. They had deliberately
withdrawn Jonathan's antirejection medicine because he no longer needed the transplant. His own
liver had -- as planned -- regenerated.

Jonathan, a 4-year-old with a shy smile and a love of dinosaurs, is among a small number of children
in the United States who have undergone a highly unusual type of transplant surgery, one that -- for
the few who are eligible -- offers a tremendous advantage: a normal life, free from antirejection
drugs, which suppress the immune system and increase the risk of infections, cancer and other
problems. Normally, transplant patients must take these powerful drugs for life.

In standard transplants, the diseased organ is completely removed and a new one put in its place.
What is different about the operation Jonathan and the other children had is that only part of the
recipient's liver is removed, and it is replaced with part of a donor's liver. At first, to prevent
rejection, the patient takes the usual drugs.

Then, doctors watch and wait. The liver has an extraordinary ability to regenerate, especially in
children, and the hope is that while the transplant is doing its job, what remains of the patient's
own liver will regenerate and start working again. The process can take a year or more; in
Jonathan's case, it took three.

If the liver does regenerate and grow large enough, doctors begin to withdraw the antirejection
medicines. The patient's immune system reactivates and, in most cases, gradually destroys the
transplant, which is no longer needed. Life goes back to normal, free from a daily schedule of pills
and their risks and expense.

``I think we need to promote this idea,'' said Dr.~Tomoaki Kato, Jonathan's surgeon. He works at
NewYork-Presbyterian Hospital/Columbia University Medical Center, but performed Jonathan's
transplant in 2006 at the University of Miami/Jackson Memorial Hospital.

``A lot of the transplant community is focused on how to get patients off immunosuppression, and
this is one way,'' he added.

But only a tiny fraction of transplant patients are candidates for the operation: certain children
with acute liver failure -- probably fewer than 100 a year in the United States, where 525 under 18
had liver transplants last year. The operation is a difficult one. It is longer and more risky than
a standard transplant, and surgeons caution that patients have to be selected carefully because not
all can withstand the surgery.

The surgery was first tried in Europe in the early 1990s, and later in the United States. But the
results were mixed -- the liver did not always regenerate -- and it never really caught on. (In
medical journals, it is called auxiliary partial orthotopic liver transplantation.) Dr.~Kato said
the results may have been poor because the early attempts included adults.

``I think the key is children,'' he said.

The best candidates are children with acute hepatic failure, a deadly condition in which the liver
suddenly stops working, often for unknown reasons. Although the liver might be able to recover, it
cannot do so fast enough to prevent brain damage and death from the toxins that build up. The only
way to save the life of someone with this condition is to perform a transplant -- or a partial one.
Such partial transplants do not work for chronic liver diseases that cause scarring because it
prevents the liver from regenerating.

All told, Dr.~Kato has performed the surgery on seven children, ranging in age from 8 months
(Jonathan) to 8 years, at Jackson Memorial. So far, the patient's own liver has recovered in six of
the seven children, and they no longer require antirejection drugs, Dr.~Kato said, adding that he
expected the need for the drugs to taper off soon for the seventh. In four, he described the
transplant as ``melting'' away completely on its own, but two others, including Jonathan, needed
surgery to remove a remnant or clear up an infection.

Dr.~Kato's first case was in 1997. That child spent three months in intensive care. ``We didn't
think it was successful,'' he said. But after two years, the liver had fully recovered.

``That gave us the idea this was something worth doing,'' Dr.~Kato said.

Other surgeons have tried the procedure. Dr.~Alan Langnas, director of liver transplantation at the
Nebraska Medical Center, said he had performed it on about 10 patients, mostly children, in the last
15 years. In some cases, he said, the patient's liver did not regenerate. At least one required a
second transplant.

``I think the success has always been a little mixed,'' Dr.~Langnas said. ``It depends on the
patient selection and how well their native liver recovers. But I think it is an important option
for some patients.''

Dr.~Simon Horslen, the medical director of liver and intestine transplant at Seattle Children's
Hospital, who was at the Nebraska center when the operations were done there, said: ``In the right
hands it's a wonderful technique. It is a case of those of us who have experienced it having to
convince others.''

Surgeons at Kings College in London have also performed the surgery, on 20 children ranging from 1
to 16 years old, during the last 20 years. Seventeen have survived. One needed a second transplant,
but in 14, their own livers regenerated, and, so far, 11 have been able to stop taking antirejection
drugs. In a recent article in a medical journal, the medical team from Kings College said the
operation should be considered for children who need transplants for acute liver failure.

But Dr.~J.~Michael Millis, the chief of transplantation at the University of Chicago Medical Center,
said, ``This has not been particularly successful in most of the hands that have tried it.''

He added, ``Even in Kato's series, the operative time is almost double, so the patients have to
spend almost twice as much time in the operating room, and I think that is actually the area that is
the Achilles' heel.'' (A liver transplant usually takes about six hours.)

Long operations require that patients be given large amounts of intravenous fluid, something that
children with liver failure generally cannot tolerate, Dr.~Millis said, explaining that the fluid
causes brain swelling that can kill them.

``I've been waiting for kids to do this on for a decade,'' he said. ``But by the time I get a liver
that is suitable, they're too sick. I have to get them in and out of the operating room and back to
intensive care as quickly as possible, with minimal fluids.''

Dr.~Langnas had similar concerns. ``Sometimes these kids are so sick, they literally have hours or a
day to live,'' he said. ``Under those circumstances, we want to not take any chances.''

Jonathan Nu\~{n}ez, whose family lives in Miami, had a textbook case of acute liver failure. At 8
months, he was perfectly happy and healthy, then he suddenly turned cranky and sleepy. He cried too
much, ate too little and began vomiting. He turned yellow, and his stomach and legs swelled. The
diagnosis was acute liver failure, cause unknown. The only hope was a transplant.

At Jackson Memorial, Dr.~Kato suggested a partial transplant. Jonathan's mother, Yailin Nu\~{n}ez,
said she and her husband immediately said yes, because it offered at least a chance that Jonathan
would be able to live a normal life, without immunosuppressants.

Children with acute liver failure shoot to the top of the list, and Jonathan got a transplant one
day after being listed. He had a rocky recovery, more so than most of Dr.~Kato's patients. Severe
rejection episodes required high doses of steroids. Other complications took him in and out of the
hospital for three months.

He stabilized, but his own liver did not seem to be regenerating; at one point it even shrank.
Ms.~Nu\~{n}ez never gave up hope, but after about two years, Dr.~Kato started to doubt that
Jonathan's liver would ever recover, and he even contemplated removing it to prevent problems. Then
it began to grow.

By last September, Jonathan's liver was large enough to work on its own. He no longer needed the
transplant. Doctors began decreasing antirejection drugs, and Jonathan's immune system did the rest.
In September, the transplant had been plainly visible on his CT scan. By November, it was gone.

But the transplant atrophied so fast that one spot where it had been connected to the small
intestine did not have a chance to close properly. An abscess formed, causing fevers and making
Jonathan quite sick. He needed antibiotics and a procedure to drain the infection. Two months later,
on Jan.~28, at NewYork-Presbyterian Morgan Stanley Children's Hospital, Dr.~Kato operated to remove
the abscess completely. A few days later, Jonathan and his family flew home to Miami.

``At the end of the day, I'm so glad,'' Ms.~Nu\~{n}ez said. ``I feel so fortunate that my son's
liver regenerated. The complications have been a struggle, and not knowing what caused his liver
failure haunts me to this day. But he can live a normal life without immunosuppression. That's what
matters. There is hope out there when you're given devastating news.''

``When it works, it's cool,'' Dr.~Langnas said.

``In Seattle, they are considering it,'' Dr.~Horslen said

\section{New Scrutiny on Censorship Issues for U.S. Companies in China}

\lettrine{O}{n}\mycalendar{Mar.'10}{02} the surface, Amazon.cn resembles its global siblings,
selling everything from Harry Potter books to sex toys. But a few searches of what the Web site
proclaims is ``the world's largest Chinese online bookstore'' reveals limits to this literary
universe.

A query for Zhao Ziyang -- the former Communist Party leader who was stripped of power for
supporting the 1989 democracy demonstration in Tiananmen Square -- returns no matches. The Dalai
Lama? A list of books that portray him as a dangerous ``splittist'' or that refer to the Chinese
government's hand-picked spiritual leader of Tibetan Buddhism, who goes by the same title. A search
of the words ``censorship'' and ``China'' comes back with ``censorship'' crossed out and three
Sino-themed suggestions, including a book called ``When China Rules the World.''

Amazon's U.S.~headquarters declined to comment on the issue of censorship, saying only that the
company was obliged to abide by the laws of the countries it which it operated. It does not, for
example, sell Hitler's ``Mein Kampf'' in Germany, where the book is banned.

But Amazon and other U.S.~companies doing business in China are coming under new scrutiny from
politicians and human rights groups since Google's recent announcement that it planned to stop
adhering to government demands that it censor search results in China and perhaps would pull out of
the country. Google complained that the accounts of dissidents who use the company's Gmail service
had been infiltrated, apparently by Chinese hackers.

The No.~2 leader in the U.S.~Senate, Richard J.~Durbin, has scheduled hearings for Tuesday on the
issue of global Internet freedom. Mr.~Durbin sent letters to 30 American information and technology
companies, including Amazon, Apple, eBay and Verizon, requesting information on their business
conduct in China and human rights practices.

Representatives from Google and McAfee are among those scheduled to appear at the hearing.

The senator also plans to focus on the Global Network Initiative, a coalition of Internet companies,
human rights organizations, investors and higher education institutions. The organization's members
have agreed to a set of guiding principles on human rights, including protecting ``the freedom of
expression rights of their users when confronted with government demands, laws and regulations to
suppress freedom of expression, remove content or otherwise limit access to information and ideas.''

The Global Network Initiative, which was begun in 2008, has also created evaluation and
accountability mechanisms through which to judge companies' participation.

So far Microsoft, Yahoo and Google have signed on to the voluntary code of conduct, but whether they
will adhere to it remains uncertain. It is unclear, for example, how search engines can censor
results in China and yet uphold the initiative's principles.

``We have done business in China for more than 20 years and we intend to stay engaged, which means
our business must respect the laws of China,'' Microsoft's chief executive, Steve Ballmer, said in a
recent blog post. ``That's true for every company doing business in countries around the world: we
are all subject to local laws.''

But now that Google has so publicly thrown down the gauntlet, American companies may find such
arguments do not offer blanket protection from questions about their conduct.

In a recent speech in which she condemned restrictions on Internet freedom, Secretary of State
Hillary Clinton said ``censorship should not be in any way accepted by any company from anywhere,''
adding that American companies should take a principled stand.

``This needs to be part of our national brand,'' she said.

Rights groups hope that the Global Network Initiative will encourage companies to act more
responsibly, even as U.S.~legislation making it a crime to enable censorship in other countries has
stalled. Human Rights Watch, a member of the Initiative, says that the framework is an important
step toward shifting the debate from whether to engage in censored markets to more adept promotion
of free access to information in places like China and Iran.

Arvind Ganesan, director of the business and human rights program at Human Rights Watch, said, ``Our
ultimate goal is to see an open and unfettered Internet everywhere in the world, but that won't
succeed if a company's first option is to leave those countries where this is not the case.''

He conceded, however, that the Global Network Initiative would be meaningless if companies cited
their participation in the framework but failed to push actively for a more open Internet while
expanding operations in countries like China.

Mr.~Ganesan noted that Western companies had gotten China to reconsider some censorship efforts. He
pointed to lobbying by Western technology companies last year that helped move Beijing to abandon a
plan that would have required manufacturers to pre-install filtering software, known as Green Dam
Youth Escort, on all new computers in China.

While the government said the software was aimed at pornography, experts said it could be used for
political censorship and monitoring users' Internet activities.

Other members of the initiative cite this victory in rejecting assertions by Microsoft, Amazon and
other U.S.~corporations that doing business in China legally compels them to abet government
censorship efforts.

``Anyone who makes that claim is being dishonest,'' said Rebecca MacKinnon, a visiting fellow at
Princeton University's Center for Information Technology Policy, who helped draft the Initiative.

Ms.~MacKinnon said that companies had a range of legal options in China that allowed them to operate
in the country without aiding censorship or monitoring. Google, for instance, kept its Gmail servers
outside China and thus beyond Chinese jurisdiction. Yahoo, in contrast, left itself and others,
including a Chinese journalist, open to prosecution because it kept servers in China, she said.

Ms.~MacKinnon also said that rather than submit to informal pressure -- as Microsoft did when it
took down a blog at the behest of a bureaucrat -- companies could demand that the government follow
legal precedents.

``Once you engage, you don't have to roll over and play dead the moment you get a call from a
government official,'' she said. ``There is respect for law and then there is complicity in
lawlessness.''

For now, the landscape remains murky, leaving even U.S.~customers with questions about how companies
do business in China.

Maochun Yu, a professor of East Asia and military history at the U.S.~Naval Academy, recently
ordered three books about China's counter-revolutionary purges in the 1950s from Amazon's Chinese
Web site, Amazon.cn, which is called Joyo Amazon. But a day after his \$64 purchase was scheduled to
arrive, Mr.~Yu discovered that all traces of his order had disappeared from the site, although his
credit card had been charged. Efforts to resolve the problem with Amazon's customer service center
failed.

Amazon's U.S.~headquarters did not respond to requests for comment. Other customers, both in China
and abroad, have reported no problems buying politically delicate literature from Amazon.cn.

But the episode left Mr.~Yu wondering whether someone in China had disapproved of his historical
interests and barred his order.

``Operationally speaking, the site is either very awkward or very controlled,'' said Mr.~Yu, who is
the author of two books about China, both of which are for sale on Amazon.com. ``I guess it's a case
of one Amazon, two systems.''

\section{A New Food Scandal in China}

\lettrine{A}{}\mycalendar{Mar.'10}{03} widening nationwide toxic food scandal that has its roots in
this tropical island resort area in the South China Sea has set off a rare case of public sniping
between officials.

The scandal, which has been reported by official news organizations in the last few days, is
centered on the cowpea — an ancient species of legume that includes the black-eyed pea.

Since late February, batches of cowpeas from the lush Sanya area of the island of Hainan have tested
positive for a highly toxic pesticide, isocarbophos, that is banned from use on fruits and
vegetables, according to a report on Tuesday in China Daily, the official English-language
newspaper. Tainted cowpeas from here have been found in the provinces of Hubei, Guangdong, Anhui and
Jiangsu.

The pesticide was banned in Hainan in 2004 but can still be found in remote parts of the island,
China Daily reported. Some farmers still use it because it is much cheaper than legal pesticides.
Though Hainan is well known among Chinese for its sweeping beaches and five-star resorts, it also
has a mountainous interior with rich farmland, much of it on terraced hillsides.

The outrage over the Hainan cowpeas, the latest in a series of Chinese food safety scandals in
recent years, erupted on Feb.~21, when the agriculture bureau of Wuhan, the capital of Hubei,
announced that it had destroyed 3.5 tons of toxic cowpeas from Hainan. An urgent nationwide warning
was issued by the central government's Ministry of Agriculture, and within days, cowpeas tainted
with the banned pesticide were discovered in the three other provinces.

Officials here in Sanya have criticized the Wuhan officials for breaking an ``unspoken rule'' that
officials in different cities and provinces report problems to one another rather than telling the
public, China Daily reported.

The release of the information by Wuhan officials ``did not save face for Sanya, nor did it save
face for the Ministry of Agriculture,'' Zhou Qingchong, an official in the Sanya agriculture bureau,
told China National Radio, an official news organization, according to the China Daily report.

Mr.~Zhou said Wuhan officials could have told Sanya officials about the cowpeas privately, and Sanya
would have sent out investigators.

``Wuhan is really not enough of a friend,'' he said.

Many provinces in China banned sales of Hainan cowpeas after the first tainted batches were
discovered. In the area of Dongguan, in Guangdong Province, the authorities destroyed 7.5 tons of
cowpeas on Sunday. The price of the crop has plummeted.

Because of its tropical climate, Hainan grows and exports vegetables that cannot be grown elsewhere
in China during the winter. But officials in Sanya, the island's main beach resort area, recently
decided to adopt policies to push tourism ahead of agriculture, said Du Liyin, director of the Sanya
Tourism Development Commission, in an interview on Tuesday.

Yet, Ms.~Du said, ``we don't plan to be 100 percent reliant on the tourism industry.''

In 2008, a widespread food scare gripped China when news organizations reported that dairy products
from nearly two dozen manufacturers were tainted with melamine, a toxic industrial chemical. At
least six children died and 300,000 people fell ill from the products, the government said. Senior
Chinese officials eventually ordered news organizations to stop reporting on the scandal, fearing
that the news highlighted endemic corruption within various levels of the government. Some parents
are still trying to sue for compensation.

\section{For 13th Time, Critic of China's Government Is Barred From Leaving Country}

\lettrine{C}{hinese}\mycalendar{Mar.'10}{03} security agents in Sichuan Province detained Liao Yiwu,
a prominent author and critic of the government, as he prepared to fly Monday to a literary festival
in Germany, human rights activists said.

It was the 13th time Mr.~Liao had been prevented from leaving the country. The Associated Press
reported that he had been placed under house arrest after being questioned by security agents for
four hours.

``How can this happen?'' The A.P. quoted him as saying. ``It's a cultural event, nothing political.
Such drama!''

Telephone calls on Tuesday to Mr.~Liao's home in rural Chengdu produced a recording saying that the
line was temporarily unavailable. Calls to his cellphone went unanswered.

Mr.~Liao was removed from a plane at Chengdu's airport as he prepared to fly to Germany to attend
lit.Cologne, one of Europe's largest literary festivals, where he was to read from one of his books,
``Miss Hello and the Farm Emperor: Chinese Society From the Bottom.''

``The reason for inviting Mr.~Liao was simple: he's a great writer,'' Traudle Berger, a spokeswoman
at the Cologne Festival, said in an interview on Tuesday. ``And China should be proud of such a
great writer.''

Ms.~Berger said Mr.~Liao's scheduled reading would still take place, with an actor assuming his
role. Proceeds from the ticketed event will be donated to the human rights group Amnesty
International, she said.

Last September, Mr.~Liao was barred from traveling to Berlin to attend an event affiliated with the
Frankfurt Book Fair, at which China was designated the honored guest.

A poet, screenwriter and new-journalism author, Mr.~Liao, 51, is one of China's best known and most
outspoken writers. Many of his works tell stories of people who have been left behind in the
nation's rush to economic and political prominence, characters that include prostitutes, a grave
robber, and a lavatory attendant.

His 2008 book ``The Corpse Walker,'' another view of Chinese society's lower rungs, was published to
international acclaim. His works are banned in China, but he has gained a large underground
following, and pirated versions of his works can be found in some Chinese bookstores.

Mr.~Liao was imprisoned for four years in the early 1990s after writing an epic poem, ``Massacre,''
which denounced the Chinese government's suppression of the 1989 Tiananmen Square protests. In
December 2007, when he traveled to Beijing to receive an award from the Independent Chinese PEN
Center, a writers' rights organization, he was detained by the police and sent back to Chengdu.

In a text-message exchange last month, Mr.~Liao said he had repeatedly met with Chengdu security
officials to negotiate for permission to attend the Cologne event, but was told that he had been
blacklisted by Beijing officials and forbidden to travel abroad.

In a Monday interview with the German network Deutsche Welle, Mr.~Liao said he was seated on the
plane at Chengdu's airport on Monday morning when a flight attendant approached and told him that
``someone is looking for you.''

``I asked who it was, and she said it would be best if I got my luggage,'' the newspaper quoted him
as saying. ``I got my bags, and while I was walking to the cabin door, I saw a police officer.''

Mr.~Liao said the police told him, ``You cannot continue doing whatever you want.''

``I told them there will be many readers at the festival,'' he said. ``I would like to go and meet
them and read some of my own pieces and play the traditional Chinese mouth organ, the xiao. I said
it was purely a literature festival and nothing political. They said they understood and were only
doing their job following orders from the top.''

On Monday, the PEN American Center, which like the Chinese organization is one of 145 affiliates of
the International PEN Center, called on China's president, Hu Jintao, to lift restrictions on
Mr.~Liao and other writers.

``It is hard to figure what the Chinese government hopes to accomplish by preventing one of its most
compelling literary voices from meeting with international colleagues and readers,'' Larry Siems,
who directs the American center's Freedom to Write program, said in a written statement.

Human Rights in China, a group based in Hong Kong, later published an open letter from Mr.~Liao to
German readers. In it, Mr.~Liao recounts learning to play the xiao while in prison, under the
tutelage of an old monk.

``How many other sages like my master are there among the Chinese people now?'' he wrote. ``I do not
know. How many innocent political prisoners are still imprisoned? I also do not know. But writers
like me from the bottom of society still have to write, record and broadcast, even to the dismay of
the Communist Party of China. I have the responsibility to make you understand that the life of the
Chinese spirit is longer than the totalitarian government.''

Germany's Foreign Ministry issued a statement expressing regret at China's detention of Mr.~Liao.
China's Foreign Ministry spokesman, Qin Gang, told journalists on Tuesday that Chinese citizens were
free to travel according to Chinese laws.

``We hope the relevant government will respect the law enforcement carried out by the Chinese
authorities,'' he said.

\section{Gold Is Not Only Standard for Olympians' Advertising Success}

\lettrine{F}{or}\mycalendar{Mar.'10}{03} Lindsey Vonn, Apolo Ohno and Shaun White, winning medals at
the Winter Olympics in Vancouver is their ticket to future endorsements. Their success, good looks
and compelling stories should help them add to their growing stable of sponsorship deals with a
variety of companies, marketing experts say.

But the vast majority of Olympic athletes -- even those who won medals -- will have a harder time
cashing in on their appearances at the Games. Winter sports like speedskating and bobsledding will
all but disappear from the airwaves and, for many Americans, interest in them will not grow again
until the next Winter Olympics in 2014.

\myphoto{Gold.Is.Not.Only.Standard.for.Olympians.Advertising.Success.jpg}{Left to right, Steven
  Holcomb's gold medal in four-man bobsled may be too obscure; the figure skater Mirai Nagasu put
  herself in good position for the 2014 Games; Julia Mancuso and Johnny Spillane could capitalize on
  silver medals.}

Olympic athletes also have to compete with well-established superstars in more popular sports like
Derek Jeter, LeBron James and Peyton Manning. Companies also tend to favor athletes who win gold
medals over those who win silver and bronze medals, regardless of the athlete's back story.

Critically, the scandals involving Tiger Woods and Michael Phelps have made companies skittish about
endorsing individual athletes, even successful Olympians. Some companies are sponsoring teams or
leagues instead, while others are spending less on sports altogether. Everyone is thinking twice
before signing deals, especially in this weak economy.

``After Tiger Woods, there's been a lot of restraint right now among marketers,'' said Bill Duffy,
the chairman of BDA Sports Management, which represents professional basketball players like Steve
Nash and Yao Ming. ``Honestly, many winter Olympians are not prominent enough to go mainstream.''

Still, landing a multiyear deal from a national brand is not the only way for an athlete to cash in
on his or her fame. Some athletes give motivational speeches and write books. Others endorse less
visible products like vitamins, sports equipment or medical causes. Athletes in events like the
snowboarding halfpipe may win deals from video-game makers.

Here is a partial list of the athletes who marketing experts say could win endorsement deals and who
may resurface before 2014.

Steven Holcomb: The pilot of USA-1's four-man bobsled crew who helped the United States win a gold
medal in the event for the first time in 62 years. Burly and affable, Holcomb looks as if he has
spent some time on a barstool, which makes him a strong candidate for a beer commercial, marketers
say. He has an inspiring story, having nearly gone blind after battling keratoconus, a degenerative
eye disease. The downside, marketers say, is that he was part of a four-man team and bobsledding is
an obscure sport.

Evan Lysacek: The first American male figure skater to win a gold medal since Brian Boitano in
1988. Lysacek, 24, edged out Russia's Yevgeny Plushenko, the defending Olympic champion, who later
took shots at Lysacek, a story line that marketers could mine. He also has sex appeal, said Bob
Dorfman, the executive creative director at Baker Street Advertising, which makes him a candidate
for advertisements including colognes, shampoos and ``any product served over ice.''

Julia Mancuso: She won two silver medals in Alpine skiing, yet lived in the shadow of Vonn. She
created a stir when she said the intense interest in Vonn was hurting the team. The drama that
followed -- Is Julia friends with Lindsey? -- thrust her into the spotlight and, marketers say,
could help her land a deal with a national brand.

J.R. Celski: He surprised many observers by winning a bronze medal in short-track speedskating
after slicing open his leg late last year. Celski had the advantage of racing alongside Ohno in
several events, increasing his exposure. Alas, with speedskaters, ``it's Apolo all the time,'' said
Matt Delzell, a group account director at Davie Brown Entertainment, which represents companies in
celebrity endorsement negotiations.

Rachael Flatt and Mirai Nagasu: Neither teenager won a medal in figure skating, but both planted
the seeds for future exposure at the 2014 Winter Games. Flatt, in particular, won kudos for doing
her homework during the Olympics, music to every teacher's ears.

Johnny Spillane: The American won two individual silver medals in the Nordic combined, which
includes ski jumping and cross-country skiing, including the first medal for the United States in
the event. Given the sport's obscurity, he may not win an endorsement deal from a national brand.
But he could garner interest from smaller companies. ``If an athlete wins a medal in a less
significant event like Nordic combined, it's going to be harder to get on Leno or Regis,'' said
Peter Carlisle, the managing director for Olympics and Action Sports at Octagon. ``But there are
regional and local deals. It's about reasonable expectations.''

\section{Olympian Crosby, Meet Fans of N.H.L. Crosby}

\lettrine{S}{idney}\mycalendar{Mar.'10}{03} Crosby will find out quickly if being an ``adopted son''
of Pittsburgh means he is forgiven.

Just two days after scoring the gold-medal-winning goal for Canada against the United States, Crosby
will take the ice Tuesday night for the Penguins against the Buffalo Sabres.

The game will not produce near the passion of the Olympic final, but Buffalo and Pittsburgh, the
defending Stanley Cup champion, are battling for home-ice advantage in the playoffs over the final
20-plus games of the season.

\myphoto{Olympian.Crosby.Meet.Fans.of.N.H.L..Crosby.jpg}{Canada's Sidney Crosby, right, battling
  Brooks Orpik — his Pittsburgh Penguins teammate — in the Olympic tournament.}

Crosby's team expects nothing less than a rousing ovation for him and the rest of the Penguins'
Olympians: Sergei Gonchar and Evgeni Malkin of Russia, Brooks Orpik of the United States, and
Marc-Andre Fleury of Canada.

``Aww, they'll be happy; it was Sid,'' said Penguins wing Max Talbot, who was the star of the
Penguins' Game 7 victory in the Stanley Cup finals last year. Referring to the Sharks' Dany Heatley,
another Canadian Olympian, he added: ``It's not like if Heatley had scored; no one would have cared.
They would just have been mad the United States lost.''

The gold medal game was such a transcendent contest that Penguins Coach Dan Bylsma said he even
expected ovations for the Sabres' five Olympians, including goaltender Ryan Miller, the tournament's
most valuable player. Between the two teams' various players, 7 of the 12 teams that played in the
Olympics will be represented here Tuesday.

``I'm not going to be surprised when they give a standing O to Sid,'' Bylsma said. ``And I'm not
going to be surprised when they give a standing O to Brooks Orpik. And I would almost expect them to
give one to Ryan Miller. I would if I was an American fan in the stands tomorrow night. You cheer
for Ryan Miller and you hopefully score on him.''

Crosby and the Penguins will not have that chance. Despite the possibility of more high drama
between Crosby and Miller, whom Crosby beat in overtime for the winner, the Sabres have decided to
give Miller some more rest and start Patrick Lalime.

The Penguins may have their own goalie to wonder about. Fleury was the third-string goalie the last
two weeks for Canada, but did not play.

``I'm sure he got a lot of work in practice,'' Penguins wing Bill Guerin said. ``You take a lot of
shots in that position. I'm sure he'll be ready.''

Crosby and Orpik were not expected to return to Pittsburgh until late Monday, but they are expected
to play Tuesday despite having opposed each other in Sunday's game.

``For Brooks Orpik and Sidney Crosby to be on the same flight back here today and one's wearing
gold, one's wearing silver and they'll be playing tomorrow, it is kind of a unique situation,''
Bylsma said. ``Both those guys are going to be putting on a Pittsburgh Penguins jersey and meaning
it tomorrow night -- I'm confident of that.''

Miller and all other non-Canadians may be coming back to the N.H.L. regular season a little
depressed for not winning the gold, but there is a bright side. Ray Shero, the Penguins' general
manager and an adviser to the United States team, noted that since Russia was knocked out in the
quarterfinals, Gonchar and Malkin did not have the opportunity to be worn out.

``To me, I'm sure they're disappointed, but they got four or five more days of rest,'' Shero said,
``so as the manager of the team, that's important too.''

But Gonchar, who had won silver and bronze medals for Russian in previous Games, the emotions will
be hard to forget.

``It's my fourth Olympics,'' said Gonchar, who wants to make his country's team in 2014, when the
Olympics will be played in Sochi, Russia. ``This time I was not going for the experience; I had
before. So now everything was about getting a gold medal, and we didn't get it, so it's probably one
of the biggest disappointments so far in my career.''

On the other side, the Penguins hope to ride the emotional high of Crosby, their captain, to etch
their name yet again on the Stanley Cup. ``It's a great thing for him and great for his confidence,
another thing on his r\'esum\'e,'' Shero said, ``as if he needs something else, right?''

Though Pittsburgh fans can be brutal to opposing teams and players, Guerin, who played on the United
States team that won the silver medal at the 2002 Games and in numerous other international
competitions, said he thought local loyalties would transcend national allegiances when it came to
Crosby.

``You know what?'' Guerin said. ``He's Pittsburgh's adopted son, and I think everyone should be
extremely proud of Sidney.''

\section{Chinese Economy Treads Risky Path. (Ask Japan.)}

\lettrine{T}{he}\mycalendar{May.'10}{26} White House dispatched 200 dignitaries to China's capital
for relationship-cementing talks this week, a number that bespeaks a world order rapidly turning
upside down. From China's threefold economic growth in a decade to today's swarms of five-star
hotels and six-figure autos, it is a clich\'e to say that the world has never before seen anything
like this juggernaut.

Actually, the world has. Twenty years ago, Japan was the juggernaut, a silky economic machine
poised, according to conventional wisdom and more than a few best sellers, to dominate world trade
and global diplomacy in the decades to come. When a flu-ridden President George Bush passed out at a
Tokyo state dinner in 1992, vomiting on the leg of the Japanese prime minister, it was widely seen
as emblematic of a recession-sapped West supplicating before a rising East.

That it never happened hardly foreshadows China's future prospects. Yet as outsiders behold China's
transformation from peasant nation to economic colossus, the risks of extrapolating from China's
robust present into an indeterminate future are not to be ignored. China's rise is doubtless an
economic miracle. But like Japan's, China's state-driven economic model may prove more useful in
earlier stages of development than in guaranteeing sustained growth well into the future.

``Japan had an economic model that worked phenomenally well for, what, 40 years? And then it stopped
working,'' Arthur Kroeber, the managing director of Dragonomics, a Beijing-based economic
forecasting firm, said in an interview. ``Does that prove the model for the first 40 years was
wrong? No -- it proves that it was right for that stage of Japan's development.''

Many of the complaints that United States officials brought to this week's Strategic and Economic
Dialogue in Beijing -- about markets closed to foreign competition; about unfairly appropriating
hard-won American technologies; about rigged currency values that swell American trade deficits --
may ring familiar to those who followed United States-Japan relations two decades ago.

For much of the postwar 20th century, Japan built an empire of private companies closely tied to and
favored by the state. It erected a trade colossus whose exports were pumped up -- American critics
said -- by an artificially depressed yen. Japanese manufacturers were accused of arrogating American
technologies to churn out low-cost electronics. Japan's retail and financial markets were all but
impenetrable to American competitors.

The United States trade deficit with Japan outraged Congress and helped prod Washington in the 1980s
to orchestrate a wholesale appreciation of the yen's value against the dollar to help protect
American manufacturers.

China today also has a stable of mega-corporations, although unlike Japan's they are explicitly
state-owned and often viewed as instruments of government policy. Broad sectors of the economy,
including finance, communications, energy and some crucial manufacturing sectors, are effectively
off limits to foreigners and even to most domestic competitors.

The renminbi, China's currency, is in global opinion -- save in China -- held unduly low to keep
China's export machine revved up. When Treasury Secretary Timothy F.~Geithner arrived here on
Sunday, an opening argument to his counterparts was that Chinese rules were forcing American
companies to surrender their technological jewels just for a ticket to compete in the Chinese
market.

Some of these tactics come straight from the playbook of Japan and other developing nations that
have sought to raise their export-driven economies to a higher level.

China is different in two respects that may seem contradictory. On one hand, major industries like
oil, telecommunications, banking and aviation are deemed strategic and are under tight state
control. Of the 22 Chinese corporations listed on the Fortune Global 500, 21 are controlled by
China's central government or state-run banks. Just one, Shanghai Automobile, is run by a local
government. None are privately owned.

These ``national champions,'' as the government deems them, are the vanguard of China's push into
global markets, and the evangelists of Chinese economic values.

On the other hand, light industry, retailing and the nation's booming export sector are more free to
play by Adam Smith's rules. In contrast to Japan, in China Western retailers and consumer goods,
from Wal-Mart to Snickers to Tesco, are ubiquitous and compete vigorously with homegrown
competitors. And many of China's leading exports, like iPods and Nike sneakers, are manufactured by
or for foreign multinationals that retain most of the profits from their sale.

By Chinese logic, there is no contradiction. Like free speech, human rights and any number of other
societal basics, private enterprise in China is vibrant and open to outsiders, but only as long as
it does not threaten the state's interests.

If this formula seems anathema to Western capitalists, Beijing's three decades of jackrabbit growth
speak for themselves. Even a senior Treasury Department official in Beijing for this week's meetings
said that China's fiscal stewards had basically perfected the art of balancing conflicting forces in
an economy of 1.3 billion people.

Less clear, many economists allow, is how long this formula will work. China itself is rapidly
changing: already, low-cost competitors like Vietnam are siphoning some of the labor-intensive
industries that powered Beijing's rise. The worldwide economic slowdown and the fact that China
already dominates some crucial industries suggest that it can no longer count on ever-rising exports
as a major source of growth.

Beijing has pledged to ease its reliance on exports by raising domestic consumption. But that
requires a turnabout in spending habits by a population accustomed to saving for contingencies, like
education and medical care, rather than spending. Turning that around may take time.

And then there is China's debt. Some influential economists argue that China has grown in part by
seizing people's savings to finance high-speed trains, steel factories and speculative real-estate
investments.

How productive that state-led investment will turn out to be is a matter of debate. It has clearly
aided China's development up to now. But just as exports cannot increase endlessly, at some point
the returns from building new roads and factories are likely to diminish as well.

The central government has surmounted several debt crises in the past decade, bailing out its banks
after unwise investments. Many analysts anticipate more bailouts after the government flooded the
economy with cheap money in the wake of the 2008 financial crisis.

Japan's economic miracle ended with the collapse of a phenomenal real-estate bubble and was worsened
by a series of policy errors that led to a long period of stagnation. China's economy is still at an
earlier stage of development, and its government is likely to seek to avoid the mistakes that
exacerbated Japan's slump.

But whether the state-driven investment and export-led development that has made China an economic
power can be replaced soon by new sources of growth has become a more pressing question.

``The Chinese fundamentals are good,'' Huang Yasheng, a professor at the Massachusetts Institute of
Technology's Sloan School of Management, said in an interview. ``All I am saying is that the policy
has not changed. You cannot rely on an artificial external stimulus to keep your economy going.''

\section{At Harvard, Kagan Aimed Sights Higher}

\lettrine{O}{ne}\mycalendar{May.'10}{26} Saturday afternoon in March 2003, Lawrence H.~Summers
invited Elena Kagan for a private chat in the library of Elmwood, the stately clapboard mansion he
occupied as the president of Harvard. The two had been close colleagues, if not close friends, as
top aides working for President Bill Clinton. But this was no social call.

Mr.~Summers, a brilliant but impatient economist with a bull-in-the-china-shop management style, was
looking for a new law school dean. Ms.~Kagan, a newly tenured professor, was thin on management
experience, and her academic writings were relatively scant.

But she was a faculty favorite -- her colleagues viewed her as a leader and consensus builder who
held sway with the strong-willed university president. Mr.~Summers, aware that she had her sights
set even higher, accompanied his job offer with a hint of a warning.

``I would say Elena's colleagues chose her as much as I did,'' he said in a recent interview,
adding, ``I said to her: 'Elena, if you accept this job and then you are offered a position like
Supreme Court justice or attorney general, I will congratulate you with all my heart and wish you
well. But we need you to make a commitment to the law school for a few years before taking any other
position.' ''

Now Mr.~Summers is President Obama's top economic adviser, and Ms.~Kagan is the president's Supreme
Court nominee. Her dealings with Mr.~Summers -- she persuaded him to abandon an unpopular plan to
move the law school, kept her distance when he faltered and made no bones about trying to succeed
him when he was forced to resign as Harvard's president -- reveal a woman of intense ambition and
deft political skills. Their relationship hints as well at Ms.~Kagan's persuasiveness, and how she
might operate on a divided Supreme Court, where persuasion often seems in short supply.

``He is not someone you can cajole in any way,'' Martha L.~Minow, the current Harvard Law dean, said
of Mr.~Summers. ``It's the merits, evidence, substance. It's not about charm, it's not about small
talk, it's 'Just the facts ma'am,' and build your case and be unbelievably fair-minded about the
other case, because he is going to ask you about the other side.''

In the days since Mr.~Obama nominated her, much has been made of Ms.~Kagan's tenure as dean, and her
top-to-bottom transformation of Harvard Law School. Far less attention has been paid to how she
climbed from visiting professor to possible successor to Mr.~Summers, the result of relentless
networking and a remarkable ability to navigate the treacherous waters of Harvard's internal
politics.

``She is a strategic and deliberative thinker on all issues -- there was always a sense of 'Let me
do my homework,' '' said one colleague, Prof. Charles J.~Ogletree Jr.~He recalled how, as dean,
Ms.~Kagan met with every member of the faculty. ``She was willing to work seven days a week, it
wasn't just Monday through Friday, 10 to 5, it was whenever people were available -- a baseball
game, a student reception, a breakfast, a lunch, a coffee.''

Ms.~Kagan had been a close ally of Mr.~Summers, but she was noticeably silent when he ran into a
public relations buzz saw over his impolitic remarks about women's aptitude for science -- remarks
that helped cost him his job. Privately, Ms.~Kagan told friends and colleagues that she thought the
fracas was overblown. But she also resisted entreaties by allies of Mr.~Summers to publicly defend
him, according to two people familiar with the situation, speaking on condition of anonymity.

One said friends of Mr.~Summers viewed her silence as an act of disloyalty, the other said Ms.~Kagan
simply did not want to drag the law school into the fray. Her style was one of careful balance, said
Dennis F.~Thompson, a professor of political philosophy.

``She doesn't avoid getting engaged in issues that are controversial,'' he said. ``But she doesn't
herself want to be the object of controversy.''

As a Harvard law graduate, Ms.~Kagan was no stranger when she arrived in Cambridge, Mass., as a
visiting professor in 1999, the year Mr.~Summers became Mr.~Clinton's Treasury secretary. The job
was ``clearly a look-see'' said Carol Steiker, a Harvard law professor and close friend of
Ms.~Kagan's, with the understanding that she would get tenure if it worked out.

Professor Kagan was an instant hit with students, demanding and energetic, with a self-deprecating
wit. She threw herself into the rhythms of faculty life, attending workshops to comment on
colleagues' writings, advising the law review board, having frequent dinners with colleagues.

Among the most contentious issues then facing Harvard was how to expand beyond the confines of
Cambridge. The university had been buying up property across the Charles River in the Boston
neighborhood of Allston. When Mr.~Summers's predecessor, Neil L.~Rudenstine, suggested that the law
school move, the law faculty responded by voting against it, 37 to 1.

``It's rare that anyone here agrees about anything,'' Professor Steiker said, ``but everyone agreed
we didn't want our campus moved across the river.''

But when Mr.~Summers was installed as president in 2001, he put the Allston move back on the table.
With the faculty up in arms, the longtime law school dean, Robert C.~Clark, asked Ms.~Kagan to head
a study committee; he thought she had the potential to succeed him as dean and wanted to give her a
leadership role. Ms.~Kagan, who had just been granted tenure, would be wading into perhaps the most
fractious issue at the university.

She seemed to have an instinctive feel for how to build a case that would work with Mr.~Summers.
``Her approach was to give a rational basis, instead of just an emotional one, for the faculty's
reaction,'' Professor Clark said.

She hired a consultant, and persuaded the university to foot the bill, producing a 101-page
strategic plan that considered everything from future growth to dormitory space to the intellectual
benefits of remaining near the arts and sciences buildings. She made no explicit recommendation, but
the study strongly suggested that Allston was far better suited to the biomedical sciences than the
law school. The plan was soon dead, and Ms.~Kagan gained folk hero status.

``I didn't think we had a snowball's chance,'' said Prof. J.~Mark Ramseyer, a member of the
committee. ``Her stock went way up on the reputational grapevine on the strength of that report,''
said another law professor, Randall L.~Kennedy.

The report came out in November 2002, the same month that Mr.~Clark announced he was stepping down
as dean. Mr.~Summers appointed a search committee to advise him on a replacement. Ms.~Kagan had
already written an award-winning article in her area of scholarship, administrative law. While
Mr.~Summers had ``real respect'' for her, he also had reservations, one person familiar with the
search said.

``He wasn't entirely sure he could trust her to make the right kind of scholarly judgments,'' said
the person, speaking on condition of anonymity to discuss internal matters.

While Mr.~Summers interviewed other candidates, including Robert C.~Post, the dean of the Yale Law
School, Ms.~Kagan, flew off to the University of Texas at Austin to interview for a dean's position.
When Harvard law professors met in small groups for lunch and dinner to discuss the qualities they
would like to see in a dean, Professor Ogletree said, she never said that she was interested in the
job. ``Not once,'' he said. ``Not even a hint.''

But in the end, Ms.~Kagan became Mr.~Summers's choice. He said she had impressed her colleagues --
and him -- ``by the way in which she had reached out and solicited many views, and her
consensus-building style.''

In an interview with The New York Times last year, Ms.~Kagan sounded grateful to Mr.~Summers.

``Quite a number of us who are women and relished working with Larry, and who felt that he had, in
some sense, recognized our talents, even when others wouldn't have, felt that he was given a bum rap
and were not shy about telling people so,'' she said.

They got on well: their one seeming disagreement -- over her decision to briefly bar military
recruiters from law school facilities -- was really ``a set piece,'' one colleague said, with
Mr.~Summers, who did not want to jeopardize Harvard's federal funding, working with Ms.~Kagan to
balance the university's interests against the law faculty's opposition to the military's ``don't
ask, don't tell'' policy.

Ms.~Kagan did something unusual for a law dean: she forged relationships elsewhere at Harvard,
especially with the Faculty of Arts and Sciences, the traditional seat of power. ``She made an
effort to reach out and build intellectual bridges,'' said Thomas M.~Scanlon Jr., a philosophy
professor who attended a luncheon Ms.~Kagan hosted for the arts and sciences faculty.

Those bridges would have been essential for anyone seeking higher office at Harvard. By 2007, with
Mr.~Summers gone, it was clear that Ms.~Kagan aspired to the university presidency, though
colleagues say she was careful not to appear to be openly campaigning for the job.

At the law school, professors and students believed she was a finalist for the president's job. ``We
all thought she would get it, and we were scared that she would leave and that all the changes she
had implemented would no longer be kept up,'' said Sarah Isgur, a 2008 graduate, who ran the Harvard
Law chapter of the Federalist Society, a conservative legal group.

But the search committee had heard reports that Ms.~Kagan had been harsh with her administrative
staff, and its members felt her scholarly interests were ``too narrow gauge,'' said one person
familiar with the search, speaking on condition of anonymity to discuss private deliberations. And
it might have been difficult for a law dean to ascend to the presidency; Derek Bok had done so, but
Harvard tended to choose humanists or scientists for the job.

When the historian Drew Gilpin Faust got the nod, law school students threw Ms.~Kagan a party.
Several hundred of them turned up in ``I ♥ EK'' T-shirts; The Harvard Crimson reported that
Ms.~Kagan teared up at the sight.

``Sometimes, you win by losing,'' the newspaper quoted her as saying, with her voice breaking
slightly as she addressed the crowd. ``All of you have made me feel like a real winner today.''

It was February 2007 -- the same month that a freshman Democratic senator from Illinois named Barack
Obama announced plans to run for the White House. Professor Steiker tried to console her friend by
telling her that there was a good chance a Democrat would win in 2008, and there might be bigger
things in store.

``I think she was disappointed not to be selected, as anyone would be,'' Professor Steiker said. ``I
remember saying to her, 'I can understand why you would want this, but what if a Democrat wins the
White House? You could be on the Supreme Court.' She nodded, but it wasn't like she said, 'Yeah, I'm
really glad I didn't get that.' ''

\section{Clinton and Geithner Face Hurdles in China Talks}

\lettrine{C}{hina}\mycalendar{May.'10}{26} and the United States opened three days of high-level
meetings here on Monday meant to broaden and deepen the ties between the world's largest developed
and developing economies.

But the opening session instead laid bare a recurring theme between Beijing and Washington: the
United States came with a long wish list for China on both economic and security issues, while China
mostly wants to be left alone to pursue policies that are turning it into an economic superpower.

President Hu Jintao, welcoming the 200-strong American delegation in the Great Hall of the People,
praised the ``mutually beneficial and win-win cooperation'' between the United States and China.
Such coordination, he said, had helped speed the recovery from the 2008 financial crisis.

On the crucial issue of China revaluing its currency -- something the Obama administration had
pushed for -- Mr.~Hu repeated China's past promises to make its effectively fixed exchange rate
respond more to the market, but the fact that the country's top leader mentioned reform at all
suggested it is on the leadership's agenda.

Still, Mr.~Hu also repeated that Beijing would move ``under the principle of independent
decision-making, controllability, and gradual progress.'' Translation: China alone will determine
the timing of any such move.

Economists said the deepening debt crisis in Greece, which came up immediately in the discussions on
Monday, would make Beijing more reluctant to allow its currency to appreciate in value in the
immediate future.

Treasury Secretary Timothy F.~Geithner did not mention China's currency in his opening remarks,
though he did allude to it in subsequent sessions. The administration has decided not to prod
Beijing at this meeting, officials said, concluding that it would resist outside pressure.

The United States is hitting similar hurdles on security issues. Secretary of State Hillary Rodham
Clinton pressed China to support measures against North Korea following the strong evidence that it
torpedoed a South Korean warship in March. But China has been skeptical of North Korea's role and is
reluctant to punish the North, with which it has close ties.

And while China agreed to a watered-down United Nations resolution on Iran's nuclear program, it has
not signed off on amendments against specific Iranian citizens and companies. With big planned
investments in Iran's oil and gas industry, China may well be in business with some of them.

In her speech to the opening session, Mrs.~Clinton cited Iran and North Korea as issues in which
Beijing and Washington must find common cause. ``Today, we face another serious challenge provoked
by the sinking of the South Korean ship,'' she said. ``So we must work together, again, to address
this challenge and advance our shared objectives of peace and stability.''

A spokesman for the Foreign Ministry, Ma Zhaoxu, was noncommittal, saying of the Korea crisis, ``We
hope all the relevant parties will exercise restraint and remain cool-headed.''

Some of this is cultural, to be sure. Chinese officials tend to speak far less directly than
Americans. Mr.~Hu did not mention Iran and North Korea, referring to regional ``hot spots.'' The
fact that he frankly addressed the exchange rate of China's currency, the renminbi, surprised some
observers, and lent itself to varying interpretations.

For some experts, Mr.~Hu's pledge to ``steadily advance the reform mechanism of the RMB exchange
rate,'' without repeating his previous references to the rate being ``basically stable,'' was a sign
of conciliation. ``It's important, the fact they haven't mentioned it,'' said Ben Simpfendorfer, the
China economist for the Royal Bank of Scotland.

But others interpreted it as a pre-emptive move to take the issue off the table. Eswar Prasad, an
economist at Cornell University, noted that the crisis in Greece had rattled the Chinese on two
levels. It was likely to curb their exports to Europe, and it had strengthened the renminbi relative
to the swooning euro, which makes Chinese goods more costly in foreign markets.

``That double hit on China's exports almost certainly means that they're not going to move forward
unless there is evidence of stabilization in the euro and stabilization in Europe's recovery,''
Mr.~Prasad said.

A senior Chinese official said that Beijing would keep a ``high alert and attention on the euro zone
sovereign debt crisis.'' He noted that it could affect not only Europe's economic recovery but also
Chinese exports. China exports more to the European Union than to the United States.

The United States needed a 48-vehicle motorcade to ferry its delegation to this second round of the
so-called strategic and economic dialogue. Among the prominent names: the chairman of the Federal
Reserve, Ben S.~Bernanke, the commander of the military's Pacific Command, Adm. Robert F.~Willard,
and the secretary of health and human services, Kathleen Sebelius.

Some of the topics under discussion veered far from economics and security. Mrs.~Clinton singled out
Melanne Verveer, the State Department's ambassador at large for women's issues, who is meeting with
Chinese women's groups to discuss their progress in women's rights.

Mr.~Geithner lobbied against Chinese government procurement rules giving preference to products with
intellectual property developed in China. American businesses, particularly in technology, say this
handicaps them and deprives China of state-of-the-art products. ``Innovation flourishes best when
markets are open, competition is fair, and strong protections exist for ideas and inventions,'' he
said.

The Chinese have their pet issues as well: Beijing is pushing Washington to loosen controls on
exports of high-technology equipment with potential military applications. A raft of questions from
reporters for state-run Chinese media organizations suggested a coordinated campaign.

If American officials seemed likely to leave China with many of their wishes unfulfilled, there was
one notable difference in this year's meeting compared to the one last year in Washington: the
American economy is growing again, which gave Mr.~Geithner a rare chance to crow a bit.

Rather than identify the United States with the troubled economies of Europe, Mr.~Geithner said the
United States was holding its own with emerging economies like Brazil, India, and China.

``Economic growth in the U.S.~and China is broader and stronger than many had anticipated, even a
few months ago,'' he said.

\section{Do You Take This Immimgrant?}

\lettrine{T}{he}\mycalendar{Jun.'10}{13} retired mechanic from Michigan looked shell-shocked beside
his bride, a classical pianist from Moscow who clutched the printed e-mail exchanges of their
Internet romance. Young newlyweds from Long Island, still recovering from their reception for 600
guests the previous weekend, faltered as their lawyer quizzed them on the details of their City Hall
ceremony four months before. A Manhattan woman bickered with her Turkish spouse about the kinds of
questions they had been warned to expect.

Did they know the color of each other's toothbrush? The pattern of the bathroom tile? What had they
done last New Year's? And were they ready to answer far more intimate queries from a government
official hunting for signs their marriage was fake?

``Embarrassing questions,'' explained the Manhattanite, Lindsay Garvy-Yeguf, 28, the butterfly
tattoo on her foot growing jittery, as her husband, Gunes Yeguf, 31, turned paler in his dark suit.
``They might ask you about your sex life.''

These three were among dozens of couples inside federal immigration headquarters in Manhattan one
recent Tuesday, seated in a crowded waiting room where posters exhorted everyone to ``Celebrate
Citizenship, Celebrate America.''

Having flunked their first interviews with United States Citizenship and Immigration Services, they
had entered the mysterious world of the ``Stokes unit,'' a uniquely New York variation on the
marriage interviews conducted nationwide whenever a citizen seeks a green card for a foreign spouse.
Named for a 1976 federal court settlement that gave couples, among other protections, the right to
bring a lawyer to a second, recorded interview if their first one raised suspicions of fraud, the
Stokes unit recently doubled its staff to 22 officers.

It is a story line familiar from pop culture: ``The Proposal'' last year, ``Green Card'' in 1990.
And while the authorities do not question the validity of the marriage of Faisal Shahzad, the failed
Times Square bomber, his arrest last month did renew questions about the process of scrutinizing
spousal green-card petitions. Nationwide, the number of such petitions denied for fraud is tiny: 506
of the 241,154 filed by citizens in the last fiscal year, or two-tenths of 1 percent (an additional
7 percent were denied on other grounds, like failing to show up for an interview).

Some critics contend that the low numbers simply show the system is easily fooled, while others say
that exaggerated estimates of marriage fraud over the years have created a bureaucratic monster,
thwarting legitimate, if unconventional, couples and spurring unconstitutional intrusion into their
lives.

In some parts of the country, the authorities stage dawn bed checks. ``Someone shows up at your
house with a badge and a gun, unannounced,'' said Laura Lichter, an immigration lawyer in Denver. ``
'Hi, we're here from immigration. Do you mind if we come in to look and see if two towels are wet?'
''

While Stokes makes such home visits off-limits in New York State, lawyers and immigrant advocates
complain that, at its worst, the process is a Kafkaesque version of ``The Newlywed Game,'' with dire
consequences: those who fail can be put on a path to deportation. Couples' futures together depend
on proving separately to a skeptical bureaucrat that, as the law states, they did not marry
``solely'' for a green card. (Passing the interviews simply makes a foreign spouse eligible for a
green card; getting one requires a separate application and security clearance.)

In the modern jambalaya of online dating, arranged weddings, bicoastal relationships, open marriages
and serial divorce, a bona fide union can be harder than ever to discern, leaving lovers who are
unable to produce a land-line telephone bill facing questions about birth control. That Tuesday in
the Stokes unit, one couple volunteered that the wife was eight weeks pregnant only to have the
husband be asked: ``Is it yours?''

``The latitude that officers have is broad, and one that has to be exercised with a lot of care,''
said Andrea Quarantillo, the immigration agency's district director for New York. ``Is it perfect?
No.~It's judgmental.''

According to an agency worksheet that was recently leaked online, red flags include ``unusual
cultural differences,'' a large age discrepancy, an ``unusual'' number of children and a citizen
with meager means. Daniel Lundy, an immigration lawyer, said the boxes on the worksheet ``pretty
much invite racial profiling and other stereotypes.''

``You could be married 50 years and still find it difficult to pass,'' Mr.~Lundy said of the Stokes
process.

The unit's lore is worthy of its own reality TV series -- sham couples caught red-handed, yes, but
also quirky ones whose authenticity surprised everyone. The gay man who claimed he had suddenly
found his female soul mate (denied); the recovering alcoholic who had lost his memory (approved);
the man who volunteered that he had erectile dysfunction in an attempt to explain why his mate did
not know the location of his nine tattoos (unsuccessful); the elderly citizen who lost an arm in a
subway accident, but found happiness with a young Caribbean wife (successful).

``We can't impose our definition of marriage, especially being in New York,'' said Maria Guerra, a
Stokes supervisor. ``We've seen it all.''

AN officer looked out on the waiting room, trying to read body language. Were some couples over the
top, snuggling and holding hands? Did some seem like strangers?

In one corner were the newlyweds from Long Island: Ersan Kahyaoglu, 25, a robustly built
electrician, American born of Turkish descent, his rumpled shirt hanging out; and Dilek Kahyaoglu,
21, Turkish but having lived in America for a decade, slim in her fitted black suit.

``I think it's impossible to know all the answers,'' Mr.~Kahyaoglu complained. ``I don't remember
what color her dress was at the civil ceremony. It was, like, different colors.''

His wife prompted him, gently: ``Black and white, flowery.''

``Who cooks?'' asked their lawyer, Raj Jadega, practicing. ``I cook,'' she said proudly.

``When was the last time you made dinner?'' She looked blank.

``Then you don't cook,'' he said.

It would have been funny if the stakes were not so grave. The couple was an automatic Stokes
referral because the wife was in deportation proceedings. She had come with her family on a tourist
visa at 11 and was 13 when their applications for green cards derailed because their first
``lawyer'' was not one.

``The questions can be arbitrary and very detailed, and they're on the firing line right now,''
Mr.~Jadega said, as the couple discussed the fine points of each other's favorite music and food
(Techno? More like rock. Chicken parmigiana or stuffed peppers? Chicken.) ``If a certain number of
questions are answered incorrectly,'' he said, ``they can stop the interview right there.''

It was hours later when the nervous groom emerged from his interview. The officer had mainly asked
``basic stuff,'' he said: ``birth dates, the type of house we live in, how much rent we pay -- maybe
25 questions.''

But then: was his wife on birth control?

``I said no,'' Mr.~Kahyaoglu said. ``He said that could mean using condoms. I said, 'No sir, we're
not using anything.' ''

Separately, his wife was pressed about condom use, and said, ``Once in a while.''

``How am I supposed to explain it to him?'' the groom asked later. `` 'Well, sometimes I feel like
reaching into the drawer by the bed -- '?''

Despite the discrepancy, they passed, and left rosy-cheeked, elated.

THE idea of marriage as a gateway for terrorists, and prosecutions of scandalous fake-marriage
rings, have periodically stirred alarm. In 1986, an estimate by the immigration agency that one in
three marriages were counterfeit spurred Congressional preoccupation and tough laws, but turned out
to be a gross exaggeration -- it was later revised to 8 percent.

An agency audit of marriage fraud, conducted in 2007, has never been released. When The New York
Times filed a request for such data under the Freedom of Information Act, the agency identified 656
relevant pages, but blacked out 655, saying the information would disclose the deliberative process
or law enforcement techniques. The Times has appealed.

Agency officials said they could not provide New York statistics to compare with national numbers,
or even count what share of New York couples were referred to the Stokes unit, because of a flawed
computer system. Officials did agree to track results during two weeks in April: 93 of the 114
couples interviewed were approved, or 81 percent. In addition to the 21 interviewed couples whose
denial was pending, 2 other marriages were deemed invalid because prior divorces were flawed, and 22
couples were no-shows denied out of hand.

On that Tuesday, 5 of the 25 couples in the Stokes waiting room were headed to rejection because of
discrepancies, though none were among those The Times had interviewed beforehand, and none were
classified as fraudulent.

``We're not mind readers,'' Ms.~Guerra, the Stokes supervisor, said. ``If we lived with them for a
month, we might see they have a marriage. But if they get everything wrong in the interview, we have
to deny the case.''

Stokes supervisors said that officers were trained to avoid questions about sex, but they cannot
stop X-rated answers, like one from the citizen wife who was asked what she did for her spouse's
birthday and began recounting their night together in explicit detail. Some couples offer
photographic evidence in the mistaken belief that the government requires proof of a marriage's
consummation.

``It's not something a normal couple would do,'' Ms.~Guerra said. ``They're overcompensating.''

Occasionally the cubicle turns into a confessional, with one spouse revealing infidelities -- even
children -- unknown to the other. And though the criminal penalties for marriage fraud are up to
five years in prison and a \$250,000 fine, a few citizens risk prosecution, withdraw their petition
and exit by a back door, leaving officers to break the news to the immigrant spouse.

Barbara Felska, a star of the Stokes unit, said her message to couples was: ``Do not fear Stokes if
your marriage is real -- all you need is love!'' But her method is to seek evidence of three aspects
of a marriage: legality (like valid divorce decrees); the commingling of assets and other joint
documentation (which young or poor couples often lack); and ``mental and emotional connection
through shared life experience.''

The last is the most subjective, and Ms.~Felska, 38, a naturalized citizen from Poland who won the
green card lottery, likes to dream up creative questions while doing dishes. Like: ``What piece of
jewelry means the most to your wife?''

On this Tuesday she was examining the three-year marriage of Yusuf Mohammed and Sally Bines.
Mr.~Mohammed, 42, a twice-divorced and Muslim taxi driver, had not seen his sons in Ghana for seven
years; his first American wife had dropped her petition for his green card. Ms.~Bines, 41, a
Christian divorc\'ee of Puerto Rican descent, was studying to be a teacher at Hostos Community
College, where they met.

They were ushered past a cubicle where an officer could be heard demanding, ``But are you the
bill-payer?'' and again, louder, after some Spanish translation, ``So you're the bill-payer, yes or
no?''

Ms.~Felska has a lighter touch. She tells couples: ``I should never know more after interviewing
your husband or your wife for 45 minutes as a federal officer than you know about her or him after
two years of marriage.'' The key moment, she finds, is when couples try to explain mismatched
answers.

Indeed, Ms.~Bines left that afternoon in a huff over one of her husband's responses. ``Some jewelry
he said he bought me that he didn't buy, he's going to buy me right now!'' she announced, as he
tried to explain that he had always meant to buy her that necklace.

The very real quarrel helped convince Ms.~Felska that the marriage was real, too.

ANOTHER union approved that day was the Internet romance between the retired mechanic, Larry
Christiansen, 66, and his Russian bride, Alla, 53, who has a Ph.D. in music and a daughter at Baruch
College.

Theirs was the kind of international courtship that Congress moved to regulate in 2005, requiring
that the citizen disclose any criminal record to the foreigner --the first such law to cast the
American husband as potential villain. It was a striking shift from 1986, when Congress passed the
Immigration Marriage Fraud Amendments after hearings that harped on ``the devious foreign husband
trying to dupe some poor American citizen into marrying him,'' said Kerry Abrams, a legal scholar at
the University of Virginia.

Mr.~Christiansen, a white-haired man with an arthritic gait who was divorced in 2001 after 35 years
of marriage, called his new union ``the best thing that ever happened in my life,'' and was shocked
when ``they even asked me if she paid me to marry her!''

The couple's lawyer, Irina Matiychenko, was pleased to report that when the officer reunited them
and asked, ``Who's the boss?'' there was instant disagreement.

``He said, 'I'm the boss.' And she said, 'No, you're not' '' -- a good sign that she was not playing
the part of docile mail-order bride vulnerable to an abusive husband, Ms.~Matiychenko explained.

Yes, there was a cultural gulf. ``I'm basically a redneck and she's very sophisticated,''
Mr.~Christiansen said. ``The princess marrying the frog.

``She's a sweetheart,'' he added. ``That's what drew me to her. When she smiles, it's real.''

NO lawyer helped Miguelina Montalvo Diaz, a 32-year-old mother from Yonkers, show that she and her
Dominican husband of six months had a life together. The letter summoning them listed documents to
bring, including a letter from their bank about joint accounts. So they quickly opened a joint
account.

Red flag! Documents dated close to the interview immediately raise suspicion. A Stokes officer asked
Ms.~Diaz to explain. ``I didn't want to have a joint account,'' Ms.~Diaz said she protested. ``You
guys asked for that!''

But the couple emerged later arm-in-arm, giddy. Their separate answers to questions like ``Where do
you keep the hamper?'' and ``Where do you keep the shoes?'' were ``100 percent on target,'' said
Ms.~Diaz, who was recently laid off from a department store, and describes her husband, Ramon Emilio
Diaz, who has four sons in the Dominican Republic, as a wonderful father to her 3-year-old boy. ``We
did good.''

Still, they were not home free. The immigration officer had noticed that Mr.~Diaz still had a
separate account, in addition to their joint one. ``They said he needs to put me on that account,''
Ms.~Diaz said, adding, ``My mom has been married 25 years and they don't have a joint account.''

She was not about to argue. In the balance was the green card that would let Mr.~Diaz easily visit
the Dominican Republic and return. And because they had been married less than two years, Mr.~Diaz
could be granted only a ``conditional'' green card. In two years, they could be back in the Stokes waiting room, facing another round of personal questions from another stranger.

\clearpage
\renewcommand\listfigurename{\textit{Table of Figures}}
{\footnotesize\textit{\listoffigures}}

\end{document}
